\documentclass[11pt]{article}
\usepackage[margin=1in]{geometry}
\usepackage{amsmath,amssymb,amsthm,booktabs}
\usepackage{hyperref}
\hypersetup{hidelinks,pdftitle={Four-zero tail-palindromic polynomials over F2: reciprocity and transformation boundaries},pdfauthor={Wencong Huang, Baohua Li}}
\newtheorem{proposition}{Proposition}[section]
\newtheorem{lemma}[proposition]{Lemma}
\newtheorem{theorem}[proposition]{Theorem}
\newtheorem{corollary}[proposition]{Corollary}
\theoremstyle{definition}
\newtheorem{conjecture}[proposition]{Conjecture}
\theoremstyle{remark}
\newtheorem{remark}[proposition]{Remark}
\title{\Large Four-zero tail-palindromic polynomials over $\mathbb F_2$:\\
reciprocity and transformation boundaries}
\author{Wencong Huang, Baohua Li}
\date{Local revision, 22 November 2026}
\begin{document}
\maketitle
\begin{abstract}
We classify the full-degree quadratic rational pullbacks that produce
monic polynomials over $\mathbb F_2$ with constant term one and exactly
four zeros in a palindromic tail. Every such output has degree $2$
modulo $4$. Cubing all roots of a four-zero candidate destroys tail
palindromicity. For fifth powers, we determine exactly which candidates
have palindromic images and give their zero counts. Sparse norm
identities and finite certificates cover all admissible degrees and
positions in these root-power classifications. For the B family, we
determine the irreducible factors with roots in the quadratic extension
of the descent field and, on an explicit parameter class, in its cubic
extension. A closed-orbit bound controls root counts over extensions
of arbitrary degree. A reciprocal gcd formula gives an exact count
of the position pairs that pass the sieve. In degree $2p$ with $p$ an odd
prime, root ratios and products give criteria for irreducibility and
primitivity. We enumerate $2663$ primitive candidates through degree
$200$ and verify a primitive example at every even degree from $6$
to $300$.
\end{abstract}

\section{The family and the question}\label{sec:intro}
Fix an even degree and consider a monic binary polynomial with
constant term one and a palindromic tail containing exactly four zeros.
The zeros occur in two reflected pairs. We determine which polynomials
of this form arise from quadratic rational pullbacks and when taking
powers of the roots preserves tail palindromicity.

Put $n=2k\ge6$ and $S_n(x)=1+x+\cdots+x^n$. We consider
\begin{equation}\label{eq:candidate}
 f_{n,i,j}(x)=S_n(x)+x^i+x^j+x^{n-1-i}+x^{n-1-j},
 \qquad 1\le i<j\le k-1,
\end{equation}
over $\mathbb F_2$. These are exactly the monic degree-$n$ polynomials
with constant term one whose nonleading coefficient vector is a
palindrome with four zero entries. The restrictions on $i,j$ exclude
the constant position. There are $\binom{k-1}{2}$ candidates.
Let $N_4(n)$ denote the number of primitive polynomials in this family.
Here primitive means irreducible with a root of multiplicative order
$2^n-1$; standard background is in \cite{LN}.

\begin{conjecture}[$C_4$]\label{conj:C4}
For every even $n\ge6$, one has $N_4(n)>0$.
\end{conjecture}

Primitive polynomials with prescribed coefficients are studied systematically
by Cohen \cite{Coh06}. In the present family the field is fixed, all but four
tail coefficients are one, and the four missing positions are coupled by
reversal. The number of coefficient constraints therefore grows with the
degree. A tail palindrome can also be encoded by the self-reciprocal
polynomial $xf(x)+1$ of odd degree $n+1$. This describes its support;
irreducibility and primitivity refer throughout to $f$ itself.

Quadratic rational pullbacks provide a classical construction of
irreducible polynomials of doubled degree \cite{MP17}.
For the four-zero family, coupled translation identities and binary
submask arguments give complete lists of the translated $Q$ and
translation-invariant outputs
(Theorems~\ref{thm:quadratic} and~\ref{thm:translation-classification}).
Combining the quadratic branches gives the degree restriction
$n\equiv2\pmod4$ and excludes iterated doubling from an even input
degree (Corollary~\ref{cor:all-quadratic}). The trace obstruction
also gives uniform reducibility for translation-invariant outputs
when the number of missing tail coefficients is divisible by four
(Proposition~\ref{prop:translation-zero-obstruction}).

The root-power transform preserves degree by recording each
powered root with its multiplicity. Cubing destroys tail
palindromicity, and fifth powers
give a complete list of palindromic images whose zero counts are
$2$ modulo $4$ (Theorem~\ref{thm:root-powers}). Sparse norm
identities reduce this classification to affine exponent collisions.
Appendix~\ref{sec:root-certificates} gives the certificates that
cover the degree and position parameters.

The B family in the quadratic classification has an explicit
descent. Propositions~\ref{prop:B-quadratic-extension}--\ref{prop:B-cubic-extension}
identify factors detected in the quadratic and cubic extensions
on their stated parameter ranges and bound the number of roots
in extensions of arbitrary degree.

Section~\ref{sec:reciprocal-screen} counts the position pairs for
which the reciprocal gcd is trivial.
Proposition~\ref{prop:root-pairs} gives root-ratio and root-product
tests in degree twice an odd prime. Section~\ref{sec:computations}
reports the complete primitive counts and first witnesses, and
Section~\ref{sec:limits} states the remaining existence questions.

\section{The reciprocal obstruction}
\label{sec:reciprocal-screen}

Tail symmetry turns polynomial reciprocity into an explicit common-factor
test. In the four-zero family, the forced factor is determined by two
integer gcds, and the number of surviving position pairs has a closed
formula. The calculation uses the binomial gcd identity \cite{LN} and
unit counts expressed by the second Schemmel function \cite[Section~1]{Def20}.

\begin{lemma}[Reciprocity and multiplicities]
\label{lem:reciprocal-multiplicity}
Let \(f=x^n+\sum_{r=0}^{n-1}c_rx^r\), where \(n\) is even,
\(c_0=1\), and \(c_r=c_{n-1-r}\). Write
\(f^*(x)=x^nf(1/x)\). Then
\[
 f^*=xf+x^{n+1}+1,\qquad
 \gcd(f,f^*)=\gcd(f,x^{n+1}+1).
\]
This gcd is squarefree. In particular, every self-reciprocal
irreducible factor of \(f\) occurs with multiplicity one. If a
non-self-reciprocal irreducible \(g\) and \(g^*\) both divide \(f\),
the smaller of their multiplicities is one.
\end{lemma}
\begin{proof}
The coefficient symmetry gives the displayed identity, hence the gcd
identity by subtraction. The derivative of \(x^{n+1}+1\) is \(x^n\),
which is coprime to it; thus its divisors are squarefree. A
self-reciprocal factor has the same multiplicity in \(f\) and \(f^*\).
For a reciprocal pair the two multiplicities are interchanged, so
their minimum is the multiplicity in the gcd. This proves the claims.
\end{proof}

\begin{proposition}[Exact reciprocal common divisor]
\label{prop:reciprocal-gcd}
For \(f=f_{n,i,j}\) from \eqref{eq:candidate}, put
\[
 N=n+1,\quad \Delta=j-i,\quad R=n-1-i-j,\quad
 a=\gcd(N,\Delta),\quad b=\gcd(N,R),\quad h=\gcd(a,b).
\]
Then, as an exact polynomial quotient in \(\mathbb F_2[x]\),
\begin{equation}\label{eq:reciprocal-gcd}
 \gcd(f,f^*)=
 \frac{(x^a+1)(x^b+1)}{(x+1)(x^h+1)},\qquad
 \deg\gcd(f,f^*)=a+b-h-1.
\end{equation}
Consequently, irreducibility of \(f\) requires \(a=b=1\).
\end{proposition}
\begin{proof}
The nonleading terms of \(f\) occur in pairs, so \(f(1)=1\).
By Lemma~\ref{lem:reciprocal-multiplicity}, with
\(S_n=(x^N+1)/(x+1)\),
\[
 \gcd(f,f^*)=\gcd(f,S_n)=\gcd(S_n,P),\qquad
 P=x^i(1+x^\Delta)(1+x^R).
\]
The oddness of \(N\) makes \(S_n\) squarefree. For positive $u,v$,
the Euclidean algorithm on the exponents gives
$\gcd(x^u+1,x^v+1)=x^{\gcd(u,v)}+1$ over $\mathbb F_2$.
Applying this identity yields
\[
 \gcd(S_n,1+x^\Delta)=\frac{x^a+1}{x+1},\qquad
 \gcd(S_n,1+x^R)=\frac{x^b+1}{x+1}.
\]
Since \(S_n\) is squarefree, its gcd with their product is their
least common multiple. Their intersection is
\((x^h+1)/(x+1)\), proving \eqref{eq:reciprocal-gcd}.
If this gcd were all of \(f\), monicity and equal degrees would give
\(f=f^*\), hence \((x+1)f=x^N+1\) and \(f=S_n\), contrary to its
four zero coefficients. Thus a nonconstant gcd is a proper divisor.
The right side is constant exactly when \(a=b=1\).
\end{proof}

\begin{corollary}[An explicit infinite reducible family]
\label{cor:infinite-reducible}
For every odd \(m\ge3\), the valid position pair
\((n,i,j)=(3m-1,m-2,m)\) satisfies
\begin{equation}\label{eq:infinite-reducible}
\begin{split}
 f_{3m-1,m-2,m}(x)
 ={}&(1+x+\cdots+x^{m-1})\\
 &\quad\cdot(x^{2m}+x^{m+1}+x^{m-1}+x^{m-2}+1).
\end{split}
\end{equation}
In particular it is reducible, and the displayed forced divisor
has unbounded degree \(m-1\).
\end{corollary}
\begin{proof}
The inequalities \(1\le m-2<m\le(3m-3)/2\) make the positions valid.
Here \(N=3m\), \(\Delta=2\), and \(R=m\), so
\(a=1\), \(b=m\) in Proposition~\ref{prop:reciprocal-gcd}.
For the explicit quotient, group the geometric sum as
\(S_{3m-1}=S_{m-1}(1+x^m+x^{2m})\). The four deleted terms sum to
\[
 x^{m-2}(1+x^2)(1+x^m)=S_{m-1}x^{m-2}(x+1)^3.
\]
Adding these expressions and cancelling the two \(x^m\) terms
inside the second factor proves \eqref{eq:infinite-reducible}.
\end{proof}

\begin{proposition}[Number of positions surviving the screen]
\label{prop:reciprocal-survivors}
For odd \(N=n+1\), let \(\varphi\) be Euler's totient and put
\[
 \mathfrak S_2(N)=N\prod_{p\mid N}\left(1-\frac2p\right).
\]
The number of four-zero position pairs with \(\gcd(f,f^*)=1\) is
\begin{equation}\label{eq:screen-survivors}
 E_N=\binom{\varphi(N)/2}{2}-\frac{\mathfrak S_2(N)-1}{2}.
\end{equation}
Thus \(N_4(n)\le E_{n+1}\). This is only an upper bound;
\(E_{n+1}>0\) does not imply \(N_4(n)>0\).
\end{proposition}
\begin{proof}
Set \(u=i+1\), \(v=j+1\). The two coprimality conditions in
Proposition~\ref{prop:reciprocal-gcd} are precisely that \(u-v\)
and \(u+v\) be units modulo \(N\). Since \(2\) is invertible,
the map \((u,v)\mapsto(u-v,u+v)\) is a bijection on
\((\mathbb Z/N\mathbb Z)^2\), giving \(\varphi(N)^2\) ordered
pairs before excluding zero coordinates. The cases \(u=0\) and
\(v=0\) contribute \(\varphi(N)\) each and are disjoint.
The independent sign changes of \(u,v\) and their interchange
act freely on the remaining pairs: admissibility excludes
\(u=\pm v\), and a nonzero coordinate cannot equal its negative.
Dividing by eight gives
\(\varphi(N)(\varphi(N)-2)/8\) unordered pairs of distinct
nonzero sign classes.

Choose representatives \(1,\ldots,k\) of those classes. The
allowed positions use \(2,\ldots,k\), so the class represented
by \(1\) must be removed. The Chinese remainder theorem counts
\(\mathfrak S_2(N)\) residues \(v\) for which both \(v-1\) and
\(v+1\) are units: modulo each prime \(p\mid N\) two residues
are forbidden. One residue is \(v=0\); the others form pairs
\(\{v,-v\}\). Hence the removed class has
\((\mathfrak S_2(N)-1)/2\) neighbors. Subtraction proves
\eqref{eq:screen-survivors}.
\end{proof}

\begin{remark}[Factorization examples and fixed positions]
\label{rem:screen-scope}
For example,
\[
 f_{8,1,2}=(x^3+x^2+1)(x^5+x^2+1)
\]
has trivial reciprocal gcd and is reducible. Repeated factors also occur:
\[
 f_{14,2,5}=(x^2+x+1)(x^6+x^4+x^3+x+1)^2.
\]
The squarefreeness hypothesis in the classical
Stickelberger--Swan factor-parity criterion \cite{Swan62} must
therefore be checked. An odd number of factors still need not mean
irreducibility: \(f_{12,4,5}\) has three distinct factors of degrees
\(3,4,5\).

For fixed \(i,j\), one also has
\[
 (x+1)f_{n,i,j}=x^{n-1-j}A_{i,j}+B_{i,j},
\]
where \(A_{i,j}=x^{j+2}+(x+1)(x^{j-i}+1)\) and
\(B_{i,j}=1+(x+1)(x^i+x^j)\) are fixed polynomials.
This places fixed-position factor-parity questions in the
classical framework of Fredricksen, Hales and Sweet \cite{FHS86}.
\end{remark}


\section{Quadratic-transform intersections}
\label{sec:quadratic}
\subsection{The translated \texorpdfstring{$Q$}{Q} branch}
For a degree-$m$ polynomial $h$, define
\begin{equation}\label{eq:quadratic-transform}
 \mathcal T(h)(x)=(x+1)^m h\!\left(\frac{x^2}{x+1}\right),\qquad
 \mathcal M_n(f)(x)=(x+1)^n f\!\left(\frac{x}{x+1}\right).
\end{equation}
The map $\mathcal T$ is the classical $Q$-transform followed by the
translation $x\mapsto x+1$; explicitly, $\mathcal T(h)(x)=h^Q(x+1)$.
The equivalence between
$\mathcal M_{2m}(f)=f$ and $f=\mathcal T(h)$ is a special case of
Mattarei and Pizzato \cite[Lemma 9]{MP17}: take $a=b=1,c=0$,
numerator $x^2$ and denominator $x+1$ in their notation.
The nondegeneracy conditions hold, and the representation is unique;
indeed the summands $x^{2j}(x+1)^{m-j}$ have distinct degrees $m+j$.
More general invariant representations are given by Reis
\cite[Theorem 1.3]{Reis20}. We combine these invariant representations
with the four-zero coefficient condition to determine every compatible
output, including reducible polynomials.

Write $\Delta a=a(x+1)+a(x)$, and use reversal with respect to a
specified formal degree when the leading coefficient may vanish.

\begin{lemma}[A descent for tail palindromes]\label{lem:descent}
Let $f=x^{2k}+\sum_{j=0}^{2k-1}c_jx^j$, with $c_0=1$ and
$c_j=c_{2k-1-j}$. Set $d=k-1$ and
\[
 e=\sum_{a=0}^{k-1}c_{2a}x^a,\qquad o=x^d e(1/x).
\]
Then
\begin{equation}\label{eq:descent}
 \mathcal M_{2k}(f)=f
 \quad\Longleftrightarrow\quad
 \Delta e=0\ \hbox{and}\ e=\Delta(xo).
\end{equation}
If the tail of $f$ has $2s$ zeros, the formal coefficient vector of
$e$ has $s$ zeros.
\end{lemma}
\begin{proof}
Tail symmetry gives $f=x^{2k}+e^2+xo^2$. Put
$p=xf+x^{2k+1}=xe^2+x^2o^2$, and homogenize $p$ to formal degree
$2k+1$, obtaining $P(X,Y)$. The tail condition is exactly
$P(X,Y)=P(Y,X)$. Moreover, $\mathcal M_{2k}(f)=f$ is equivalent
to $P(X,X+Y)=P(X,Y)$. Conjugating this substitution by the
interchange of $X,Y$ gives $(X,Y)\mapsto(X+Y,Y)$. Hence, for
this symmetric $P$, the condition is equivalent to $\Delta p=0$.
Direct expansion gives
\[
 \Delta p=x(\Delta e)^2+
 \bigl(e(x+1)+\Delta(xo)\bigr)^2.
\]
The two summands have disjoint odd and even supports. Thus both
vanish exactly under the conditions in \eqref{eq:descent}.
Each tail mirror pair contains one even and one odd position,
which proves the assertion about zeros.
\end{proof}

\begin{lemma}[Two missing coefficients and translation]
\label{lem:two-zero-periodic}
Let $d$ be positive and even, and $1\le a<b\le d-1$. Then
$e=S_d+x^a+x^b$ satisfies $\Delta e=0$ exactly in the following cases:
\begin{equation}\label{eq:two-zero-types}
\begin{array}{c|c|c}
 &d&\{a,b\}\\ \hline
\mathrm{I}&2^r-2,\ r\ge3&\{2^u,2^v\},\quad0\le u<v\le r-1\\
\mathrm{II}&2^r,\ r\ge2&\{2^r-1,2^t\},\quad0\le t\le r-1\\
\mathrm{III}&2^r+2^s-2,\ r>s\ge2&\{2^r-1,2^s-1\}.
\end{array}
\end{equation}
\end{lemma}
\begin{proof}
A binary submask of an integer is obtained by retaining any subset of
its nonzero binary digits; it is proper when it differs from the integer.
The binary expansion identity
$(x+1)^v=\prod_{j:\,v_j=1}(x^{2^j}+1)$ shows that
$\Delta x^v$, for $v>0$, consists of the proper binary submasks
of $v$, and has $2^{\operatorname{wt}_2(v)}-1$ terms.
Consequently $\Delta x^v=1$ exactly for powers of two.
For distinct positive $v,w$, equality $\Delta x^v=\Delta x^w$
also occurs exactly when both are powers of two: with at least
two binary digits, the singleton submasks recover all set digits.
We will also use
\begin{equation}\label{eq:geometric-translation}
 x(x+1)\Delta S_t=x^{t+2}+(x+1)^{t+2}+1.
\end{equation}
It follows that $\Delta S_t=0$ exactly when $t+2$ is a power of two.

For $0\le b\le t$, the one-missing-coefficient polynomial
$S_t+x^b$ is translation invariant exactly when
\[
 (t,b)=(2^s-2,0)\quad\hbox{or}\quad(2^s-1,2^s-1),\qquad s\ge1.
\]
Indeed equality of its values at $0,1$ forces $b=0$ if $t$ is even.
If $t$ is odd, translation invariance forces its leading term to
cancel, since a nonzero invariant polynomial has even degree, so
$b=t$. Equation~\eqref{eq:geometric-translation} completes both cases,
including the zero polynomial at $t=b=0$.

Put $d=2u$ and write $e=E^2+xO^2$.
Odd and even coefficient comparison gives
$\Delta e=0\Longleftrightarrow\Delta O=0,\ \Delta E=O$.
The useful residual is
\[
 B_u=S_{u-1}+\Delta S_u
     =\frac{x^{u+1}+(x+1)^{u+1}+1}{x}.
\]
If $a,b$ are even, then $O=S_{u-1}$. Thus $u+1$ is a power
of two and $B_u=0$. The remaining condition is
$\Delta(x^{a/2}+x^{b/2})=0$, so $a/2,b/2$ are powers of two.
These are precisely the even-even pairs in type I.

If $a,b$ are odd, then $E=S_u$ and
$B_u=x^{(a-1)/2}+x^{(b-1)/2}$. Its two terms force
$\operatorname{wt}_2(u+1)=2$. Writing $u+1=2^v+2^w$, $v>w\ge0$,
gives $B_u=x^{2^v-1}+x^{2^w-1}$. This yields type III when $w>0$;
$w=0$ yields type II with $t=0$.

Finally suppose the missing positions are $2A$ and $2B+1$.
Then $O=S_{u-1}+x^B$ and $B_u=x^B+\Delta x^A$.
The one-missing-coefficient classification gives either
$u=2^s-1,B=0$, or $u=2^s,B=u-1$.
In the first case $B_u=0$; in the second $B_u=x^{u-1}+1$.
Both force $\Delta x^A=1$, so $A$ is a power of two.
They give respectively type I with one missing position $1$,
and type II with $t>0$. These cases exhaust the parities.
All displayed pairs satisfy the required equations on substitution,
so the classification is sufficient as well as necessary.
\end{proof}

\begin{theorem}[Complete four-zero intersection]\label{thm:quadratic}
For the candidates \eqref{eq:candidate},
$\mathcal M_n(f_{n,i,j})=f_{n,i,j}$ if and only if one of the
following holds:
\begin{align}
 n&=2^R+2^S-2,\quad R>S\ge2,&
 (i,j)&=(2^S-2,2^S-1);\label{eq:quadratic-A}\\
 n&=2^R+2,\quad R\ge3,&
 (i,j)&=(3,2^{R-1});\label{eq:quadratic-B}\\
 n&=14,&&
 (i,j)\in\{(2,4),(2,5),(4,5)\}.\label{eq:quadratic-C}
\end{align}
Equivalently, these are exactly the four-zero polynomials in the
image of $\mathcal T$.
\end{theorem}
\begin{proof}
We first use coefficient comparison to reduce the four-zero condition to
the two-zero polynomial in Lemma~\ref{lem:descent}. We then apply the
three types in Lemma~\ref{lem:two-zero-periodic} and impose the second
descent identity on each type.

Write $n=2k$. The coefficient of $x^{n-1}$ in
$\mathcal M_n(f)$ is $f'(1)=k-2$ in $\mathbb F_2$, since the
four missing tail positions include two odd positions: each pair
$\{r,n-1-r\}$ contains one odd position because $n-1$ is odd.
Its value in $f$ is $1$, so $k$ is odd. Comparison of the
$x^2$ coefficients gives $[x]f=\binom{2k}{2}=k=1$.
Thus positions $1,n-2$ are not missing. In Lemma~\ref{lem:descent},
$e=S_d+x^a+x^b$ is therefore monic of even degree $d=k-1$,
with $1\le a<b\le d-1$.

By Lemma~\ref{lem:two-zero-periodic}, $e$ has one of the three types
in \eqref{eq:two-zero-types}. The second descent condition is
equivalent to
\begin{equation}\label{eq:second-descent}
\begin{split}
 x^a+x^b+\Delta x^{d+1-a}+\Delta x^{d+1-b}
 =\frac{x^{d+2}+(x+1)^{d+2}+1}{x}.
\end{split}
\end{equation}
For type I the right side is zero. The binary masks $d+1-a$ and
$d+1-b$ have all $r$ digits set except the distinct digits specified
by $a,b$. Their proper-submask supports have symmetric difference
of size $2^{r-1}-2$: among all submasks the difference has size
$2^{r-1}$, and deleting the two top masks removes two terms.
Equation~\eqref{eq:second-descent} requires that difference to be
the two terms $x^a+x^b$, so $r=3$. For $r=3$ it is exactly that
pair, giving $d=6$ and the three cases \eqref{eq:quadratic-C}.

For type II put $\{a,b\}=\{d-1,2^t\}$. The equation becomes
\[
 \Delta x^{d+1-2^t}=x^{2^t}+x+1.
\]
For $t=0$ both sides are $1$. For $t\ge1$ the right side has
three terms, while $d+1-2^t$ has $r-t+1$ set digits. Thus
$2^{r-t+1}-1=3$, forcing $t=r-1$, which also suffices.
The two choices yield respectively \eqref{eq:quadratic-A} with
$S=2$, and \eqref{eq:quadratic-B}.

For type III the two $\Delta$ terms are both $1$ and cancel.
The right side is $x^a+x^b$, so every such pair satisfies the
condition. Restoring the even zero positions $2a,2b$ and their
tail mirrors gives \eqref{eq:quadratic-A} with $R=r+1,S=s+1\ge3$.
All these steps are reversible, proving the assertion in every degree.
The equivalent image statement follows from the classical
representation preceding Lemma~\ref{lem:descent}.
\end{proof}

\begin{corollary}[Reducibility and failure of iteration]
\label{cor:quadratic-limits}
Every member of \eqref{eq:quadratic-A} is reducible.
If $h$ is itself a four-zero polynomial of the form
\eqref{eq:candidate}, then $\mathcal T(h)$ is never a member of
that family. Among the three polynomials in
\eqref{eq:quadratic-C}, only $f_{14,2,4}$ is irreducible.
\end{corollary}
\begin{proof}
For $f=\mathcal T(h)$ let $m=n/2$; its constant term gives $h(0)=1$,
and the $x^2$ coefficient gives
\begin{equation}\label{eq:quadratic-trace}
 [x]h=[x^2]f+\binom m2.
\end{equation}
In \eqref{eq:quadratic-A}, $S=2$ gives $m\equiv1\pmod4$ and
$[x^2]f=0$, while $S\ge3$ gives $m\equiv3\pmod4$ and $[x^2]f=1$.
Hence $[x]h=0$ in either case. If $h$ is reducible, multiplicativity
of the transform makes $f$ reducible. If $h$ is irreducible, the
classical binary $Q$ criterion \cite[Theorem 1.1]{Ugo13} says that
$h^Q$, and hence $f=h^Q(x+1)$, is irreducible exactly when $[x]h=1$.
This proves the first assertion.

The theorem, or its coefficient proof, gives $n\equiv2\pmod4$
for every four-zero output. A four-zero input has even degree,
so its transformed degree is divisible by four. This proves the
iteration obstruction for this particular transform.

For the exceptional pairs $(2,4),(2,5),(4,5)$ the corresponding
$h$ are, respectively,
\[
 x^7+x^6+x^3+x+1,\quad
 x^7+x^6+x^3+x^2+x+1,\quad
 x^7+x^6+x^4+x^2+1.
\]
The first satisfies $x^{128}\equiv x\pmod h$ and
$\gcd(h,x^2+x)=1$, so is irreducible by the prime-degree Rabin test;
its linear coefficient is $1$. The second has root $1$,
and the third has zero linear coefficient. The same $Q$ criterion
finishes the proof.
\end{proof}

\begin{proposition}[Explicit descent for the B family]\label{prop:B-descent}
For $R\ge3$, put $q=2^{R-1}$. The unique polynomial with
$\mathcal T(h_R)=f_{2q+2,3,q}$ is
\begin{equation}\label{eq:B-descent}
 h_R(x)=1+x+x^2+(1+x)\left(x^q+
       \sum_{t=2}^{R-2}x^{q-2^t+1}\right),
\end{equation}
where the sum is empty for $R=3$. Moreover, $\deg h_R=q+1$ and
$[x]h_R=1$.
\end{proposition}
\begin{proof}
Set $L(x)=\sum_{a=0}^{R-2}x^{2^a}$ and
\[
 F(x)=x^{q+2}h_R(1/x)
     =(x^2+x+1)x^q+(x+1)(L(x)+x^2).
\]
For $s=X^{-1}$ and $w=s+s^2$, telescoping gives $L(w)=s+s^q$;
substitution and collection of terms give
\[
 s^{-2q-4}F(w)=(X+1)f_{2q+2,3,q}(X).
\]
The left side is $(X+1)\mathcal T(h_R)(X)$, proving the formula;
the distinct leading degrees of the transform basis give uniqueness.
The exponents in the sum lie between $q/2+1$ and $q-3$, so the
degree and linear coefficient follow directly.
\end{proof}
The classical $Q$ criterion makes irreducibility of $h_R$ equivalent
to that of its output.

For the remainder of this subsection put $k=R-1$, $q=2^k$, and
$f_R=\mathcal T(h_R)=f_{2q+2,3,q}$.

\begin{proposition}[Quadratic extension and the degree-$12$ factor]
\label{prop:B-quadratic-extension}
For every $R\ge3$,
\[
 \gcd(f_R,x^{q^2}+x)=1,\qquad
 \gcd(h_R,x^{q^2}+x)=
 \begin{cases}
    x^2+x+1 & k\equiv3\pmod4,\\
    P_{12} & k\equiv6\pmod{24},\\
    1 & \text{otherwise},
 \end{cases}
\]
where $P_{12}=x^{12}+x^{11}+x^9+x^8+x^7+x^6+x^5+x+1$ is
irreducible. Over the full parameter range, $P_{12}\mid h_R$
exactly when $R\equiv7\pmod{24}$. The quadratic has multiplicity
two and $P_{12}$ has multiplicity one when present. Their pullbacks
are irreducible of degrees $4$ and $24$, with the same multiplicities.
\end{proposition}

In particular, $f_R$ has no irreducible factor of degree at most $B$
if $\operatorname{lcm}(1,\ldots,B)\mid 2(R-1)$.

\begin{proposition}[Closed Frobenius orbits]\label{prop:B-closed-orbits}
For every $R\ge3$ and $m\ge1$,
\[
 \#\{x\in\mathbb F_{q^m}:h_R(x)=0\}\le 9^m,\qquad
 \#\{x\in\mathbb F_{q^m}:f_R(x)=0\}\le 2\cdot 9^m.
\]
In particular, an irreducible factor of $h_R$ of degree $d$ satisfies
$d\le 9^{d/\gcd(d,k)}$.
\end{proposition}

\begin{proposition}[Cubic extension on the class $24\mid k$]
\label{prop:B-cubic-extension}
Suppose $24\mid k$. The two gcds are as follows.
\begin{center}
\begin{tabular}{@{}ccc@{}}
\toprule
$k\bmod 72$ & $\gcd(h_R,x^{q^3}+x)$ & $\gcd(f_R,x^{q^3}+x)$\\
\midrule
$0$  & $1$ & $1$\\
$24$ & $P_{18,1}$ & $\mathcal T(P_{18,1})$\\
$48$ & $P_{18,2}$ & $\mathcal T(P_{18,2})$\\
\bottomrule
\end{tabular}
\end{center}
The polynomials $P_{18,1}$ and $P_{18,2}$ are irreducible of
degree $18$. In hexadecimal coefficient notation, in which bit $i$
is the coefficient of $x^i$, they are \textup{\texttt{0x412b1}} and
\textup{\texttt{0x55abf}}, respectively. The first pullback splits into two
distinct irreducible factors of degree $18$. The second is irreducible
of degree $36$.
\end{proposition}

The two nontrivial rows each give an infinite progression of reducible
four-zero outputs. Appendix~\ref{sec:B-certificates} proves the three
propositions and supplies the finite certificates for the exact
factor classifications.

\subsection{The translation branch}\label{sec:translation}
We next impose $f(x+1)=f(x)$. The invariant representation
$f=H(x^2+x)$ and its Artin--Schreier irreducibility criterion are
classical \cite[Lemma 2.1 and Theorem 2.2]{Kap17}. The additional
constraints below come from the tail palindrome and its missing
coefficients.

\begin{lemma}[Degree compatibility]\label{lem:translation-degree}
If a monic tail-palindromic polynomial $f$ of even degree $n\ge6$
satisfies $f(x+1)=f(x)$, then $n+2$ is a power of two.
\end{lemma}
\begin{proof}
Let $F(X,Y)=Y^nf(X/Y)$ and $N=n+1$. Tail symmetry gives
\[
 XF(X,Y)+YF(Y,X)=X^N+Y^N.
\]
By translation invariance, put
$A=F(X,Y)=F(X+Y,Y)$, $B=F(Y,X)=F(X+Y,X)$ and
$C=F(X,X+Y)=F(Y,X+Y)$. Applying the displayed identity to
three pairs of variables yields
\[
\begin{aligned}
 XA+YB&=X^N+Y^N,\\
 (X+Y)A+YC&=(X+Y)^N+Y^N,\\
 XC+(X+Y)B&=X^N+(X+Y)^N.
\end{aligned}
\]
Multiply these equations by $X+Y$, $X$ and $Y$, respectively,
and add. All terms on the left cancel, leaving
$(X+Y)^{n+2}=X^{n+2}+Y^{n+2}$. In characteristic two this
forces $n+2$ to be a power of two: if its binary expansion has
at least two nonzero digits, a singleton proper submask gives
a nonzero mixed binomial term.
\end{proof}

\begin{proposition}[A missing-coefficient obstruction]
\label{prop:translation-zero-obstruction}
Let $f$ be monic of even degree $n=2k\ge6$, with constant term one
and a tail palindrome containing $2s$ zeros. If $f(x+1)=f(x)$,
write $f=H(x^2+x)$ with $H$ monic of degree $k$. Then
\begin{equation}\label{eq:translation-zero-trace}
 [x^{k-1}]H=s\pmod2.
\end{equation}
In particular, if the number of tail zeros is divisible by four,
then $f$ is reducible.
\end{proposition}
\begin{proof}
With $e,o$ as in Lemma~\ref{lem:descent}, tail symmetry gives
$f=(x^k+e)^2+xo^2$. Odd and even coefficient comparison gives
\begin{equation}\label{eq:translation-descent}
 \Delta f=0\quad\Longleftrightarrow\quad
 \Delta o=0,\qquad \Delta(x^k+e)=o.
\end{equation}
Indeed $\Delta f=x(\Delta o)^2+
[\Delta(x^k+e)+o(x+1)]^2$.
The formal coefficient vector of $o$ has length $k$ and $s$ zeros,
so $[x^{n-2}]f=o(0)=o(1)=k-s$ in $\mathbb F_2$.
Lemma~\ref{lem:translation-degree} gives
$k=2^{R-1}-1\equiv3\pmod4$, where $R\ge3$.
The classical coefficient comparison in $H(x^2+x)$ therefore gives
\[
 [x^{k-1}]H=[x^{2k-2}]f+\binom{k}{2}=s\pmod2.
\]
Now suppose that $s$ is even. If $H$ is reducible, so is its composition.
Otherwise let $\beta$ be a root of the irreducible $H$ in
$\mathbb F_{2^k}$. Its absolute trace is $[x^{k-1}]H=0$,
so $z^2+z=\beta$ has a solution in that same field.
Thus $f$ has a root of degree at most $k$ and is reducible.
The last step is classical: specialize \cite[Theorem~2.2(2)]{Ago77}
to $p=2$, $n=r=1$ and $a=1$ in Agou's notation, with irreducible
outer polynomial $H$; see also \cite[Theorem~2.2]{Kap17}.
The same binary specialization follows from Long
\cite[Theorems~4.1--4.2, pp.~69 and 71]{Long73}; the basic
Artin--Schreier trace criterion is classical, as noted in
\cite[p.~63]{Long73}. The associated high-coefficient restriction appears in
\cite[Theorem~2.8(3)]{Kap17}.
\end{proof}

The restriction $n\ge6$ is essential: $x^2+x+1$ is a
translation-invariant irreducible tail palindrome with no zeros.
For odd $s$, \eqref{eq:translation-zero-trace} only removes this
trace obstruction; it does not establish irreducibility of $H$.

\begin{theorem}[Complete translation intersection]
\label{thm:translation-classification}
The polynomial $f_{n,i,j}$ in \eqref{eq:candidate} satisfies
$f_{n,i,j}(x+1)=f_{n,i,j}(x)$ if and only if
\begin{equation}\label{eq:translation-four}
 (n,i,j)\in\{(14,3,4),(30,5,9),(30,5,12),(30,9,12)\}.
\end{equation}
All four polynomials are reducible.
\end{theorem}
\begin{proof}
In \eqref{eq:translation-descent}, put $d=k-1$. The polynomial
$o$ is monic of degree $d$ and has two missing coefficients.
Since $k$ is odd, $o(1)=k-2=1$; invariance gives $o(0)=1$.
Thus $o=S_d+x^a+x^b$ with $1\le a<b\le d-1$, and
$e=S_d+x^{d-a}+x^{d-b}$.
By Lemma~\ref{lem:translation-degree}, $d=2^r-2$, $r\ge2$.
The case $r=2$ has only one interior position and is impossible.
For $r\ge3$, Equation~\eqref{eq:geometric-translation} gives
$\Delta S_d=0$. Hence $\Delta x^a=\Delta x^b$, which, by the
binary-support observation in the proof of
Lemma~\ref{lem:two-zero-periodic}, forces $a,b$ to be distinct
powers of two. The second condition in
\eqref{eq:translation-descent}, using $\Delta x^{d+1}=S_d$, becomes
\begin{equation}\label{eq:translation-support}
 \Delta x^{d-a}+\Delta x^{d-b}=x^a+x^b.
\end{equation}

If $a,b$ are both even, take square roots in this equation.
The two exponents on the left are $(r-1)$-digit all-one binary
masks with different singleton digits removed. Their full
submask supports have symmetric difference of size $2^{r-2}$;
removing the two top masks leaves $2^{r-2}-2$ terms.
The right side has two terms, so $r=4$.
For three-digit masks the remaining difference is precisely the
two removed singleton digits; thus all pairs from $\{2,4,8\}$
satisfy \eqref{eq:translation-support}.

Otherwise $a=1$ and $b$ is an even power of two. The term
$\Delta x^{d-b}$ has only even exponents. The mask
$d-1=2^r-3$ has all digits except digit one set, so
$\Delta x^{d-1}$ has $2^{r-2}-1$ odd terms. The right side has
only the odd term $x$, forcing $r=3$.
Here $\Delta x^5=x^4+x+1$ and $\Delta x^4=\Delta x^2=1$;
of $b=2,4$, exactly $b=4$ satisfies the equation.

A missing position $a$ in $o$ corresponds to the two positions
$2a+1$ and $n-2-2a$ in $f$. Thus $d=6$, $\{a,b\}=\{1,4\}$
gives $(14,3,4)$; $d=14$ and the three pairs from $\{2,4,8\}$
give the degree-$30$ cases in \eqref{eq:translation-four}.
The substitutions above also prove sufficiency.
Reducibility follows uniformly from
Proposition~\ref{prop:translation-zero-obstruction} with $s=2$.
\end{proof}

\subsection{Full-degree quadratic constructions}\label{sec:quadratic-scope}
Here a full-degree quadratic pullback over $\mathbb F_2$ means
\[
 F=h^mP(g/h),\qquad \gcd(g,h)=1,\qquad
 \max(\deg g,\deg h)=2,
\]
where $\deg P=m>1$ and $\deg F=2m$.
The last condition is automatic for irreducible
$P$; in general the degree drops exactly when $h_2\ne0$ and
$P(g_2/h_2)=0$ \cite[Section 2]{MP17}.

\begin{corollary}[Boundary for full-degree quadratic pullbacks]
\label{cor:all-quadratic}
If a full-degree quadratic pullback is a four-zero tail palindrome,
then it belongs either to the lists in Theorem~\ref{thm:quadratic}
or to \eqref{eq:translation-four}, and its degree is $2$ modulo $4$.
If it is irreducible, it is either $f_{14,2,4}$ or a member of
\eqref{eq:quadratic-B}. Consequently, no sequence of quadratic
pullbacks, each doubling the degree, can reach a four-zero tail
palindrome after starting at an even degree.
\end{corollary}
\begin{proof}
Write $g=\sum_{j=0}^2g_jx^j$, $h=\sum_{j=0}^2h_jx^j$, and put
\[
 a=g_2h_1+g_1h_2,\quad b=g_0h_2+g_2h_0,\quad
 c=g_1h_0+g_0h_1.
\]
The classical quadratic framework gives $b^2+ac=1$ and, in the
separable case,
\[
 (ax+b)^{2m}F\!\left(\frac{bx+c}{ax+b}\right)=F(x)
 \qquad\text{\cite[Section 4, Lemma 9]{MP17}}.
\]
In the separable case the possible triples are exactly $(1,0,1),(0,1,1),(1,1,0)$,
corresponding to reciprocity, translation and $\mathcal M_{2m}$.
There is no nontrivial scalar normalization over $\mathbb F_2$.
In the inseparable case $g_1=h_1=0$, the output is a square;
its odd coefficients vanish, whereas a four-zero candidate has
$[x^{2m-1}]F=1$.

The reciprocal branch contains no four-zero candidate: combining
$f^*=f$ with Lemma~\ref{lem:reciprocal-multiplicity} gives
$(x+1)f=x^{n+1}+1=(x+1)S_n$, hence $f=S_n$.
The other two branches are classified by
Theorems~\ref{thm:quadratic} and~\ref{thm:translation-classification}.
The irreducible alternatives follow from their reducibility
conclusions. All listed degrees are $2$ modulo $4$.
Starting at an even degree, every subsequent doubled degree is
divisible by four, proving the last assertion.
\end{proof}

The three involutions thus connect the support symmetries to an
exhaustive list of full-degree quadratic outputs.


\section{Root powers and coefficient reflection}\label{sec:root-powers}
Power substitution and taking powers of roots are different operations.
For a monic $h\in\mathbb F_2[x]$ with constant term one, define
\begin{equation}\label{eq:root-power-definition}
 \mathcal N_\ell(h)(y)=\prod_{h(\alpha)=0}(y+\alpha^\ell),
\end{equation}
where roots are counted with multiplicity. Its degree is $\deg h$,
even when a root power has a shorter Frobenius orbit. It need not be
the minimal polynomial of a single root power.
This is the classical generalized Graeffe transform
\cite[Section 5]{BFSS06}, \cite[Section 3.1]{GHL15}:
\[
 \mathcal N_\ell(h)(y)
   =\operatorname{Res}_z(h(z),y+z^\ell)
   =\det(yI+C_h^\ell),
\]
where $C_h$ is a companion matrix of $h$.
These identities retain multiplicities, including for a reducible $h$.
In contrast,
$h(x^\ell)$ has degree $\ell\deg h$; classical order criteria for its
irreducibility are recalled in \cite[Section 2]{BM25}.

\begin{theorem}[Cubic and quintic root-power images; computer-assisted]
\label{thm:root-powers}
Let $f=f_{n,i,j}$ be a candidate from \eqref{eq:candidate}.
No irreducibility or squarefreeness assumption is imposed.
\begin{enumerate}
\item The polynomial $\mathcal N_3(f)$ is never tail-palindromic.
\item The polynomial $\mathcal N_5(f)$ is tail-palindromic if and only if
\begin{equation}\label{eq:quintic-parameters}
 (n,i,j)=(10,3,4)
 \quad\text{or}\quad
 (n,i,j)=(7d-1,2d-1,3d-1),
 \quad d\ge1,\quad d\equiv\pm1\pmod {10}.
\end{equation}
The first image has two missing tail coefficients; the second has
exactly $2d$. In particular, neither root-power operation has a
four-zero tail-palindromic image.
\end{enumerate}
\end{theorem}

The sparse norm identities reduce the classification to a finite
collection of affine exponent collisions. The certificates in
Appendix~\ref{sec:root-certificates} cover all admissible degrees
and zero positions. The supplied programs verify the support and
integer-feasibility calculations.

\subsection{A sparse reflection defect}\label{sec:sparse-norm}
For odd $\ell$ and a primitive $\ell$th root of unity $\zeta$,
\begin{equation}\label{eq:norm-product}
 \mathcal N_\ell(h)(x^\ell)
   =\prod_{s=0}^{\ell-1}h(\zeta^s x).
\end{equation}
This follows by multiplying $x^\ell+\alpha^\ell$ over the roots of $h$.
It also proves that the coefficients lie in $\mathbb F_2$.
The operation is multiplicative, commutes with reciprocity, and sends
$1+x$ to $1+y$. These are elementary norm identities; the
object-specific issue is the coefficient reflection of their outputs.

Put $N=n+1$, $\Delta=j-i$, $R=n-1-i-j$, and
\begin{equation}\label{eq:feedback}
\begin{split}
 D&=x^i(1+x^\Delta)(1+x^R),\qquad u=(1+x)D,\\
 P&=(1+x)f=1+x^N+u,\qquad P^*=1+x^N+xu .
\end{split}
\end{equation}
The identity for $P^*$ follows from
Lemma~\ref{lem:reciprocal-multiplicity}. The eight raw terms of $u$
may cancel when positions are adjacent; all such cancellations are
retained. Thus $P$ has at most ten terms, independently of $n$.

Let $g=\mathcal N_\ell(f)$, $Q_\ell=\mathcal N_\ell(P)$, and
$F_\ell(x)=Q_\ell(x^\ell)$. Then $Q_\ell=(1+y)g$, so the exact
tail-reflection test for $g$ is
\begin{equation}\label{eq:norm-reflection}
\begin{split}
 \Omega_\ell(x)
  &:=
    x^{\ell N}F_\ell(1/x)+x^\ell F_\ell(x)
      +(1+x^\ell)(1+x^{\ell N})\\
  &=(1+x^\ell)
       \bigl(g^*(x^\ell)+x^\ell g(x^\ell)+1+x^{\ell N}\bigr).
\end{split}
\end{equation}
Thus $g$ is tail-palindromic if and only if $\Omega_\ell=0$:
the bracket vanishes precisely when its nonleading coefficients
are a palindrome.
In a homogeneous expansion in $u$, the highest-degree terms cancel
between $\mathcal N_\ell(P^*)$ and $y\mathcal N_\ell(P)$.
The remaining support still has bounded size for fixed $\ell$.

\begin{lemma}[Exact support certificates]\label{lem:norm-certificates}
For $k\ge3$ and $1\le i<j<k$, the polynomial $\Omega_3$ is nonzero.
The polynomial $\Omega_5$ is zero exactly at the parameter sets in
\eqref{eq:quintic-parameters}.
\end{lemma}
\begin{proof}
Expand the fixed support of $P$ in \eqref{eq:feedback}, separately
in each residue class of $(k,i,j)$ modulo $\ell$. Use the positive
gaps $i$, $j-i$ and $k-j$ to parametrize all allowed positions.
For $\ell=3$, the supplied certificate exhibits an uncancellable
term in each of $42$ parameter regions.
For $\ell=5$, the supplied collision cover has $672$ nodes and
leaves only the degree-$10$ exception and the two residue branches
of the progression in \eqref{eq:quintic-parameters}.
The construction, exhaustive coverage, and independent exact checks
are detailed in Appendix~\ref{sec:root-certificates}.
Sufficiency for the two progression branches also follows directly
from \eqref{eq:quintic-support} below; the isolated exception is
checked explicitly. These are finite parameter certificates, with
no upper bound on the free integer parameters.
\end{proof}

\subsection{The quintic progression and its zero count}
Write $f_d=f_{7d-1,2d-1,3d-1}$, where $d\equiv\pm1\pmod {10}$, and put
\[
 \mathcal A=\{7,8,9,12,13,14,21,22,23,26,27,28\}.
\]
A compact support identity for the quintic progression is
\begin{equation}\label{eq:quintic-support}
\begin{split}
 (1+y)\mathcal N_5(f_d)(y)
   ={}&1+y^{7d}
       +\sum_{a\in\mathcal A}y^{\lfloor ad/5\rfloor}\\
     &+\sum_{b=2}^{5}\bigl(y^{bd-1}+y^{bd}\bigr).
\end{split}
\end{equation}
One obtains this identity by splitting $P$ into its five exponent
residue classes and multiplying the five factors in
\eqref{eq:norm-product}. Equivalently, substitute
$d=10t+1$ or $10t+9$ into the $26$-term norm form
\eqref{eq:quintic-norm-form}; the resulting affine exponent pairs
are verified by the independent branch checker.
This proves an identity in the free parameter $t$, including
all cancellations at $d=1$.

The support symmetry is visible in the formula. The set
$\mathcal A$ is invariant under $a\mapsto35-a$, and
\[
 \lfloor ad/5\rfloor+\lfloor(35-a)d/5\rfloor=7d-1.
\]
The terms $y^{bd-1}+y^{bd}$ are paired by $b\mapsto7-b$.
Thus the right side of \eqref{eq:quintic-support} satisfies
the required reflection test.

For $d=10t+\epsilon$, the $22$ displayed exponents are distinct
when $\epsilon=1,t\ge1$ or $\epsilon=9,t\ge0$.
Order them as $0=e_0<e_1<\cdots<e_{21}=7d$.
Division by $1+y$ makes the coefficient at degree $s$ the parity
of the number of exponents $e_r\le s$. Hence the missing coefficients
occupy the intervals $[e_1,e_2-1]$, $[e_3,e_4-1]$, and so on, and
\begin{equation}\label{eq:quintic-zero-count}
 \sum_{r=0}^{9}(e_{2r+2}-e_{2r+1})
     =20t+2\epsilon=2d.
\end{equation}
The ordered affine supports and this sum are included in the
branch certificate. At $d=1$, cancellation gives
\[
 \mathcal N_5(f_{6,1,2})=y^6+y^5+y^4+y+1,
\]
again with two missing tail coefficients.
The isolated case is
\[
 \mathcal N_5(f_{10,3,4})
   =y^{10}+y^9+y^8+y^6+y^5+y^4+y^3+y+1,
\]
whose missing positions are $2$ and $7$.

\begin{proof}[Proof of Theorem~\ref{thm:root-powers}]
Equation~\eqref{eq:norm-reflection} and
Lemma~\ref{lem:norm-certificates} give the cubic obstruction and
the complete quintic parameter list.
Equations~\eqref{eq:quintic-support}--\eqref{eq:quintic-zero-count},
together with the two explicit small images, give the asserted
zero counts. Every allowed $d$ is odd, so each quintic
tail-palindromic image has $2$ modulo $4$ missing coefficients.
It therefore cannot belong to the four-zero family.
\end{proof}

\subsection{Irreducibility and degree-preserving extraction}
\begin{corollary}\label{cor:root-extraction}
For $\ell=3$ or $5$, there are no two irreducible four-zero
candidates $h,f$ of the same degree such that $h\mid f(x^\ell)$.
If an irreducible four-zero polynomial has a root of order $m$, then
\[
 \ell\not\equiv2^r\pmod m\qquad(r\ge0,\ \ell=3,5).
\]
\end{corollary}
\begin{proof}
For a root $\beta$ of $h$, the element $\beta^\ell$ is a root
of $f$ and has a full degree-$n$ Frobenius orbit.
Thus $\mathcal N_\ell(h)=f$, contrary to
Theorem~\ref{thm:root-powers}. If $\ell\equiv2^r\pmod m$,
taking the $\ell$th power merely permutes the conjugates, giving
$\mathcal N_\ell(f)=f$ and the same contradiction.
\end{proof}

The infinite quintic shape branch itself is reducible except at $d=1$:
\begin{equation}\label{eq:quintic-factor}
 f_d=S_{d-1}(x)
       \bigl(S_6(x^d)+x^{2d-1}(1+x)(1+x^{2d})\bigr).
\end{equation}
Indeed $S_{7d-1}=S_{d-1}S_6(x^d)$ and the four deleted terms are
$x^{2d-1}(1+x^d)^3$.
For $d>1$ both factors in \eqref{eq:quintic-factor} are nonconstant.
Since $\gcd(d,5)=1$, taking fifth powers permutes the nontrivial
$d$th roots of unity, so
$\mathcal N_5(S_{d-1})=S_{d-1}$; the image is also reducible.
Consequently the only irreducible inputs with a tail-palindromic
quintic image are $f_{6,1,2}$ and $f_{10,3,4}$.
Their root orders are respectively $63$ and $341$, as checked by
irreducibility and exact-order tests in the supplied program.
In particular, if $4\mid n$, the quintic image of an irreducible
four-zero candidate is not tail-palindromic.

Taking an $\ell$th root within $\mathbb F_{2^n}$ can increase a
deficient $\ell$-primary part of an element's order.
Corollary~\ref{cor:root-extraction} shows that for $\ell=3,5$ this
degree-preserving one-step construction cannot also keep the
four-zero tail.


\section{Primitive counts and reproducible computations}\label{sec:computations}
Complete enumeration of the $161700$ candidates at the $98$ even degrees
$6\le n\le200$ gives
\[
 \sum_{\substack{6\le n\le200\\ n\text{ even}}}N_4(n)=2663,
 \qquad N_4(n)>0\ \text{at every tested degree}.
\]
A separate implementation reproduces every complete count and first pair.
At the $50$ even degrees $202\le n\le300$, checking the lexicographic
prefixes covers $5410$ further candidates and verifies the first pair
at each degree.

\begin{table}[ht]
\centering\small
\begin{tabular}{rrrr}
\toprule
$n$ & candidates & $N_4(n)$ & first $(i,j)$\\
\midrule
6 & 1 & 1 & $(1,2)$\\
8 & 3 & 1 & $(2,3)$\\
10 & 6 & 2 & $(1,2)$\\
12 & 10 & 3 & $(1,3)$\\
20 & 36 & 1 & $(4,5)$\\
40 & 171 & 11 & $(1,4)$\\
62 & 435 & 16 & $(1,17)$\\
100 & 1176 & 31 & $(1,23)$\\
126 & 1891 & 55 & $(1,2)$\\
154 & 2850 & 45 & $(1,23)$\\
194 & 4560 & 20 & $(6,65)$\\
200 & 4851 & 17 & $(1,71)$\\
\bottomrule
\end{tabular}
\caption{Selected complete counts. All $98$ rows and all successful
position pairs through degree $200$ are supplied with the verification data.}
\label{tab:counts}
\end{table}

\paragraph{Arithmetic verification and reproducibility.}
The verification code constructs coefficients directly and tests
irreducibility, followed by the standard exact-order conditions
\begin{equation}\label{eq:ordertest}
 x^{2^n-1}\equiv1\pmod f,\qquad
 x^{(2^n-1)/p}\not\equiv1\pmod f\quad(p\mid2^n-1\text{ prime}).
\end{equation}
The arithmetic tests use python-flint $0.9.0$ with FLINT $3.6.0$
\cite{Flint}.
The $148$ order factorizations correspond to the $98$ complete-count
degrees and the $50$ further first-witness degrees. Exact division leaves
residual cofactor one; all $484$ distinct prime factors, of at most $144$
bits, pass FLINT's rigorous \texttt{fmpz\_is\_prime} procedure.
A standard-library implementation independently agrees on all $220$
candidates through degree $24$.

The folder \path{verification/primitivity/} supplies complete candidate
verdicts, successful positions, exact order-test residues for the $148$
first witnesses, recovered integer factorizations and input hashes.
Historical inputs and correction records are preserved with the audit.
The data and verification programs are publicly archived at
\href{https://doi.org/10.5281/zenodo.23766722}{doi:10.5281/zenodo.23766722}
and in the
\href{https://github.com/63331849121/yijing-mathematical-structures/tree/paperF-v2026.10.29-r1/PaperF_fourzero}{fixed GitHub release}.

\paragraph{Sparse feedback from primitive witnesses.}
For a primitive witness $f$ of degree $n$, the multiple $P=(1+x)f$
from \eqref{eq:feedback} defines an $(n+1)$-stage binary shift register.
Of its at most ten terms, the leading term specifies the new bit, so
each step uses at most nine old bits.
For $f_{6,1,2}=x^6+x^5+1$, the three feedback terms give
$s_{t+7}=s_{t+5}+s_{t+1}+s_t$.
Since $f(1)=1$, the factors $1+x$ and $f$ are coprime.
The classical product decomposition gives two fixed states (all zero
and all one) and two cycles of length $2^n-1$ \cite{Elspas59}.
Bitwise complementation exchanges the two long cycles, as in Lempel's
difference lifting \cite{Lempel70}.
The sparse multiple is almost primitive with increment one
\cite[Definition~3.1 and Section~8]{BZ04}.
The verified witnesses give this construction for every even $n$
from $6$ to $300$. Conjecture~\ref{conj:C4} would extend the guarantee
within the four-zero family to every even $n\ge6$.

The \href{https://doi.org/10.5281/zenodo.23880299}{version 3.0 archive}
supplies the certificates and verification programs for the root-power
proofs. Appendix~\ref{sec:root-certificates} gives the coverage argument
and file names.

The companion programs also check the reciprocal filter, both quadratic
classifications and the descent identities; detailed ranges accompany
the archived outputs.
The four translation outputs have independently certified factor degrees
$2+12$, $12+18$, $6+24$ and $7+7+8+8$, with a repeated factor in the last.
The B-sequence factorization data cover $3\le R\le10$:
the first term is irreducible and the next seven reducible.

\section{Open questions}\label{sec:limits}
\paragraph{The B family.}
The full-degree quadratic classification leaves the B family
as the only unbounded branch in which irreducibility can occur.
The explicit descent in Proposition~\ref{prop:B-descent} and
the binary $Q$ criterion reduce the question to the irreducibility
of $h_R$. The extension-field classifications give infinite
progressions of reducible members, and the closed-orbit bound
restricts factor degrees. Is $h_R$ irreducible for infinitely
many $R$?

\paragraph{Irreducibility in degree $2p$.}
Proposition~\ref{prop:root-pairs} detects irreducibility of a
given degree-$2p$ candidate by a root ratio $y$ with
$y^{2^p+1}=1$ and $y^3\ne1$, for $p$ an odd prime.
Proposition~\ref{prop:fixed-phase-labels} bounds the labels
at a fixed phase for two choices of zero positions. Do irreducible
four-zero tail-palindromic polynomials of degree $2p$ exist for
infinitely many odd primes $p$?

\paragraph{Primitive existence.}
Conjecture~\ref{conj:C4} asks for a primitive candidate in
every even degree $n\ge6$. The computations verify such
candidates through degree $300$. In degree $2p$ with $p$ an
odd prime, Proposition~\ref{prop:root-pairs} determines
primitivity from the two orders attached to a conjugate pair.

\appendix
\section{Coefficient and finite-field descriptions}\label{sec:classical-descriptions}
The coefficient quotient organizes position pairs into classes with explicit
multiplicities; the finite-field parametrization encodes primitivity
through coprime order components.

\subsection{Coefficient quotient}\label{sec:quotient}
Let $\rho$ reverse $\mathbb F_2^k$, and set
\[
 P=\{(A,\rho A):A\in\mathbb F_2^k\},\quad
 K=\{(C,C):C=\rho C\},\quad U=P/K,\quad r=\lfloor k/2\rfloor.
\]
The map $[(A,\rho A)]\mapsto A+\rho A$ identifies $U$ with
$\operatorname{im}(I+\rho)$, of dimension $r$.
Prepending a zero to $x\in\mathbb F_2^{k-1}$ defines
$\Phi(x)=[(\widetilde x,\rho\widetilde x)]$.
The all-ones tail lies in $K$, so deleting positions $i,j$ gives
class $\Phi(e_i+e_j)$. With zero-based coordinates and a suitable basis
$z_1,\ldots,z_r$, the columns $u_i=\Phi(e_i)$ form the multiset
\begin{equation}\label{eq:columns}
 \{u_1,\ldots,u_{k-1}\}=
 \begin{cases}
 \{z_1,z_1,\ldots,z_{r-1},z_{r-1},z_r\},&k=2r,\\
 \{z_1,z_1,\ldots,z_{r-1},z_{r-1},z_r,0\},&k=2r+1.
 \end{cases}
\end{equation}
Indeed pair positions $i,k-1-i$; the boundary pair contributes the
singleton $z_r$ and an unpaired middle position contributes zero.
In particular $\Phi$ is onto.

For $s=\sum s_a z_a$, write $h=\sum_{a<r}s_a$ as an integer and
$\varepsilon=k\bmod2$. The ordinary coset weight enumerator is
\begin{equation}\label{eq:enumerator}
 A_s(z)=\sum_{\Phi(x)=s}z^{\operatorname{wt}(x)}
       =(1+z)^\varepsilon z^{s_r}(2z)^h(1+z^2)^{r-1-h}.
\end{equation}
Each repeated column allows $00,11$ for zero syndrome and $01,10$
for nonzero syndrome; the singleton and optional zero column give the
remaining factors. Thus $\mu(s)=[z^2]A_s(z)$ counts four-zero candidates
in class $s$, and the number of occupied classes is
\begin{equation}\label{eq:Dk}
 D(3)=1,\qquad D(k)=1+\binom{\lfloor(k+1)/2\rfloor}{2}\quad(k\ge4).
\end{equation}
This follows by counting the nonzero quadratic coefficients in
\eqref{eq:enumerator}; at $k=3$ the zero class is unoccupied.
For a character $\chi$ of $U$, put
$t=\#\{i:\chi(u_i)=-1\}$ and $\ell=k-1$. Then
\begin{equation}\label{eq:krawtchouk}
 \widehat\mu(\chi)=\sum_{i<j}\chi(u_i)\chi(u_j)
   =K_2(t;\ell)=\frac{(\ell-2t)^2-\ell}{2}.
\end{equation}
These are classical direct-product weight-enumerator and Krawtchouk
identities \cite[Sections 3--4]{Forney11}\cite[Sections 1.1 and 2.1]{Brouwer18}.
The different factorization behavior within a class is visible already
at degree eight:
$x^8+x^7+x^6+x+1$ is primitive, while
$x^8+x^7+x^5+x^2+1=(x^2+x+1)(x^6+x^4+x^2+x+1)$ is reducible;
their tail difference $0110\,0110$ belongs to $K$.

\subsection{Finite-field parametrization}\label{sec:dickson}
Put $q=2^k$, $\mathbb K=\mathbb F_q$, $\mathbb L=\mathbb F_{q^2}$ and
$T=\{\gamma\in\mathbb L^*:\gamma^{q+1}=1\}$.
The coprime orders $q-1,q+1$ give $\mathbb L^*=\mathbb K^*\times T$.
A primitive element is uniquely $\beta=\nu\gamma$ with
$\operatorname{ord}(\nu)=q-1$ and $\operatorname{ord}(\gamma)=q+1$.
Writing $\tau=\gamma+\gamma^{-1}\in\mathbb K$, the standard Dickson
identity $D_a(u+u^{-1})=u^a+u^{-a}$ and Frobenius pairing give
\begin{equation}\label{eq:dickson}
 F_{\nu,\gamma}(x)=\prod_{i=0}^{k-1}
 \left(x^2+\nu^{2^i}D_{2^i}(\tau)x+\nu^{2^{i+1}}\right).
\end{equation}
Here $D_0=0,D_1=t,D_{a+1}=tD_a+D_{a-1}$ in characteristic two
\cite{LN}. Conversely, the two order conditions make this product the
minimal polynomial of a primitive element. Expanding the product and
comparing with \eqref{eq:candidate} gives an exact coefficient-and-order
system in which the support and the two exact-order conditions are
imposed simultaneously.

\subsection{Root ratios and products in twice prime degree}\label{sec:root-ratio}
Fix an odd prime $p$, set $q=2^p$ and take $f=f_{2p,i,j}$ from
\eqref{eq:candidate}. We use $\mathbb K,\mathbb L,T$ from
Appendix~\ref{sec:dickson} and take roots in $\overline{\mathbb F}_2$.

\begin{lemma}\label{lem:cube-subfield}
If $f$ is irreducible and $\alpha$ is one of its roots, then
$\alpha^3\notin\mathbb K$.
\end{lemma}
\begin{proof}
Write $s_m=\sum_{f(a)=0}a^m$, counting multiplicities. If $f$ is
irreducible and $\alpha^3\in\mathbb K$, its $2p$ roots occur in
pairs $a,a^q$ with equal cubes. Thus $s_m=0$ for every integer
multiple $m$ of $3$, including negative multiples since $f(0)=1$.

To compute both signs, take $C(z)=f^*(z)$ for $s_r$ and $C(z)=f(z)$
for $s_{-r}$. Denote the corresponding sum by $w_r$ and put
$c_b=[z^b]C(z)$. The logarithmic derivative of
$C(z)=\prod_a(1+a^{\pm1}z)$ gives
\begin{equation}\label{eq:signed-sums}
 \frac{zC'(z)}{C(z)}=\sum_{r\ge1}w_r z^r,\qquad
 w_r=(r\bmod 2)c_r+\sum_{b=1}^{r-1}c_b w_{r-b}.
\end{equation}
For $p\ge5$, let $h_b$ be the indicator of $b\in\{i,j\}$. The
reflected holes of $f$ start at $2p-1-j\ge p$. Substitution in
\eqref{eq:signed-sums} at $r=3$ gives
\[
 s_3=1+h_1+h_2,\qquad s_{-3}=1+h_3+(1+h_1)h_2 .
\]
Both can vanish only for $(i,j)=(1,3)$ or $i=2$, $j\ge4$.

For $p\ge11$ the following entries give $s_m=1$ on these branches.
\[
\begin{array}{c|rrrrrrrr}
i&1&2&2&2&2&2&2&2\\
j&3&4&5&6&7&8&9&\ge10\\ \hline
m&9&-9&9&-9&9&9&-15&-9
\end{array}
\]
Each entry of absolute value $9$ follows from \eqref{eq:signed-sums}
through $r=9$, using $S_9(z)+z^i+z^j$ for the negative sign and
$S_9(z)+z^{i+1}+z^{j+1}$ for the positive sign, modulo $z^{10}$.
For $(i,j)=(2,9)$ use $S_{15}(z)+z^2+z^9$ modulo $z^{16}$.
At $p=11$ the reflected hole adds $z^{12}$ to this prefix.
Both prefixes give $s_{-15}=1$.

At $p=5$ the remaining pairs $(1,3),(2,4)$ give $s_9=1$ and
$s_{-9}=1$, respectively, by \eqref{eq:signed-sums} with the full
polynomial. At $p=7$ the pairs $(1,3),(2,4),(2,6)$ give $s_9=1$,
$s_{15}=1$, $s_9=1$, respectively. The remaining $f_{14,2,5}$ is
divisible by $x^2+x+1$, as one checks modulo this factor using $x^3=1$.
At $p=3$ the only candidate is $x^6+x^5+1$, with $s_3=1$.
This exhausts the allowed positions.
\end{proof}

\begin{proposition}[Root ratios and products]\label{prop:root-pairs}
The polynomial $f$ is irreducible if and only if two of its roots
have ratio $y\in T$ with $y^3\ne1$. It is primitive if and only if
some root ratio generates $T$ and some root product generates
$\mathbb K^*$. In the latter test both orders can be read from
the same pair $(\alpha,\alpha^q)$.
\end{proposition}
\begin{proof}
If $f$ is irreducible, $y=\alpha^q/\alpha$ belongs to $T$ and
$y^3\ne1$ by Lemma~\ref{lem:cube-subfield}. Conversely, such a $y$
has degree $2p$. It lies in neither $\mathbb K$ nor $\mathbb F_4$,
the maximal proper subfields of $\mathbb L$, since
$T\cap\mathbb K^*=\{1\}$ and $|\mathbb F_4^*|=3$ \cite[Chapter~2]{LN}.

Suppose $y=\beta/\alpha$ and $f$ is reducible. Since $y$ lies in
the compositum of the fields of the two roots, $2p$ divides the
least common multiple of their degrees. Both degrees are less than
$2p$, so one equals $p$. Exchanging the roots replaces $y$ by
$y^{-1}$ and lets us take $\alpha\in\mathbb K$ of degree $p$.
Then $\beta=y\alpha\in\mathbb L\setminus\mathbb K$. Its irreducible
factor is distinct from that of $\alpha$, has degree at most $p$,
and has degree dividing $2p$. Its degree must be $2$. As $p$ is odd,
$\mathbb F_4^*\subset T$, giving
$\alpha=\beta/y\in\mathbb K^*\cap T=\{1\}$, a contradiction.

Now let $f$ be irreducible and write $\alpha=\nu\gamma$ in
$\mathbb K^*\times T$. If $\nu=1$, then $\alpha^q=\alpha^{-1}$,
so $f=f^*$. Lemma~\ref{lem:reciprocal-multiplicity} would give
$(x+1)f=x^{2p+1}+1$, hence $f=S_{2p}$, which has no holes.
Thus $\nu$ has degree $p$. Also $\gamma^3\ne1$ by
Lemma~\ref{lem:cube-subfield}, so $\gamma$ has degree $2p$.
Every second root is $\alpha^{2^k}$ for some $0\le k<2p$.
A nonidentity ratio in $T$ forces $\nu^{2^k}=\nu$, hence $k=p$.
A product in $\mathbb K^*$ forces
$\gamma^{2^k}=\gamma^{-1}=\gamma^{2^p}$, also giving $k=p$.
For that pair the two values are
\[
 \alpha^q/\alpha=\gamma^{-2},\qquad \alpha^q\alpha=\nu^2 .
\]
These have the orders of $\gamma$ and $\nu$, since both group orders
are odd. All conjugates give the same two orders. A generator $y$ of $T$
satisfies $y^3\ne1$ since $q+1>3$, so its occurrence already gives
irreducibility. Now $\gcd(q+1,q-1)=1$ gives the primitivity assertion.
\end{proof}

Consider two central choices of zero positions,
\[
 F_{p;d,s}(x)=f_{2p,p-1-d,p-1-s}(x),\qquad
 (d,s)\in\{(6,4),(7,3)\},
\]
for a prime $p\ge11$. Fix a prime $\ell>5$ and set
$Q=2^\ell$, $D=(Q+1)/3$ and $m=2p+1$. Assume
$3(Q-1)\mid m$ and $\gcd(m,D)=1$, and use reduced residue
classes $p\bmod D$, so $\gcd(p,D)=1$.
For a root $\alpha$ of a degree-$2\ell$ irreducible factor of
$F_{p;d,s}$, write $\alpha=a u$ in $\mathbb F_{Q^2}^*$ with
$a\in\mathbb F_Q^*$ and $u^{Q+1}=1$. This decomposition is
unique since the two group orders are coprime. Retain roots
for which $S=u^3$ has order $D$. Since $a^m=1$, their phase is
$Z=\alpha^m=S^k$, where $k=m/3\bmod D$. Since $\gcd(m,D)=1$, $k$ is a unit modulo $D$ and $Z$ has order $D$.
The class $k$ determines
$p\bmod D$ by $2p\equiv3k-1\pmod D$. Its reflection label is
$\{k,-k\}$.

\begin{proposition}[Labels at a fixed phase]\label{prop:fixed-phase-labels}
Fix $\ell$ and a Frobenius orbit of phases of order $D$. As $p$
varies under the conditions above, this orbit is attained by at most
$28$ distinct reflection labels for $(d,s)=(6,4)$ and at most
$32$ for $(d,s)=(7,3)$. The same bounds hold for the signed
classes $k$.
\end{proposition}
\begin{proof}
Choose a representative $z$ of the phase orbit and take a
conjugate root $\alpha$ with $\alpha^m=z$ for each factor
being counted. The four omitted terms give
\[
 (1+x)F_{p;d,s}(x)=1+x^m+x^{p-1-d}C_{d,s}(x),\qquad
 C_{d,s}(x)=(1+x)(1+x^{d-s})(1+x^{11}).
\]
Evaluate at $\alpha$, square, and divide by $z=\alpha^m$.
Multiplication by $\alpha^{2d+3}$ then gives
\[
 H_{d,s,z}(\alpha)=0,\qquad
 H_{d,s,z}(x)=C_{d,s}(x)^2+(z+z^{-1})x^{2d+3}.
\]
For fixed $z,d,s$, the polynomial $H_{d,s,z}$ is independent of $p$.
Since $d+s+1=11$, we have $\deg C_{d,s}=2d+2$. The second term
has lower degree, so $\deg H_{d,s,z}=4(d+1)$. Each admissible
root determines $S=u^3$, which has order $D$. Thus $S^k=z$
determines at most one signed class $k\bmod D$ per root.
Identifying $k$ with $-k$ can only decrease the count.
\end{proof}

For $f_{10,3,4}$ the conjugate ratio and product have orders
$11=(q+1)/3$ and $31=q-1$, respectively. For $f_{58,15,18}$ their
orders are $q+1$ and $(q-1)/1103=486737$. The accompanying
\path{verify_root_pairs.py} reproduces both examples using \eqref{eq:ordertest}.

For a test in $\mathbb F_2[Y]$, put
\[
 R_f(Y)=\operatorname{Res}_x(f(x),f(Yx)),\quad
 P_f(Y)=\operatorname{Res}_x(f(x),x^{2p}f(Y/x)),\quad
 \Theta_p(Y)=\frac{Y^{q+1}+1}{Y^3+1}.
\]
Classical composed multiplication gives $R_f$ from $f,f^*$ and $P_f$
from $f,f$ \cite{BC87,BFSS06}. Their roots are the pairwise ratios and
products, respectively. Since $q+1$ is odd and divisible by $3$,
$\Theta_p$ has root set $T\setminus\mathbb F_4^*$ and $f$ is
irreducible exactly when $\gcd(R_f,\Theta_p)\ne1$. In that case
$\gcd(R_f,\Theta_p)$ and $\gcd(P_f,Y^{q-1}+1)$ are the minimal
polynomials of $\alpha^q/\alpha$ and $\alpha^q\alpha$, of degrees
$2p$ and $p$. Indeed, all qualifying pairs are $(a,a^q)$ by the
proof above. Their ratios and products run through single Frobenius
orbits of those lengths, and both gcds discard multiplicities.

\section{Exact parameter certificates for root powers}
\label{sec:root-certificates}
The norm identity writes the reflection defect as a sum of monomials
with affine exponents. This appendix explains how the certificates
cover the parameter regions where these exponents can collide and
how the checkers verify the arithmetic on each region. The version
3.0 archive cited in Section~\ref{sec:computations} supplies the
certificates and checkers in \path{verification/root_power/}.

\subsection{Fixed support and parameter coverage}
Represent the exponent $e_k k+e_i i+e_j j+e_0$ by
$(e_k,e_i,e_j,e_0)$. The complete input support of $P$ in
\eqref{eq:feedback}, before cancellations, is
\begin{equation}\label{eq:feedback-support}
\begin{gathered}
 (0,0,0,0),\ (2,0,0,1),\\
 (0,1,0,0),\ (0,1,0,1),\
 (0,0,1,0),\ (0,0,1,1),\\
 (2,-1,0,-1),\ (2,-1,0,0),\
 (2,0,-1,-1),\ (2,0,-1,0).
\end{gathered}
\end{equation}
In a fixed residue class $r=(k,i,j)\bmod\ell$, the coefficient
of this term in $P(\zeta^s x)$ is
\[
 \zeta^{s(e_k r_k+e_i r_i+e_j r_j+e_0)} .
\]
Multiplication, reciprocity and addition then give the full support
of $\Omega_\ell$ in \eqref{eq:norm-reflection}.
Every equality of specialized exponents is retained modulo two.

Set $a=i$, $b=j-i$, $c=k-j$.
Let $a_0,b_0,c_0\in\{1,\ldots,\ell\}$ be their positive residue
representatives and write
\[
 a=a_0+\ell A,\qquad b=b_0+\ell B,\qquad c=c_0+\ell C,
 \qquad A,B,C\in\mathbb Z_{\ge0}.
\]
This is a bijective parametrization of all admissible positions.
An exponent in \eqref{eq:feedback-support} becomes
\begin{equation}\label{eq:gap-exponent}
\begin{split}
 &\ell(e_k+e_i+e_j)A+\ell(e_k+e_j)B+\ell e_k C\\
 &\hspace{12mm}
   +e_k(a_0+b_0+c_0)+e_i a_0+e_j(a_0+b_0)+e_0 .
\end{split}
\end{equation}
A split into a variable equal to zero or at least one is disjoint
and exhaustive; in the latter branch replace it by one plus a new
nonnegative variable.

For $\ell=3$, splitting $P=p_0+p_1+p_2$ by exponent residues gives
$F_3=p_0^3+p_1^3+p_2^3+p_0p_1p_2$ by \eqref{eq:norm-product}.
Direct expansion leaves $16$--$36$ symbolic
terms per residue class. The certificate covers the $27$ residue
classes by $42$ regions: $41$ are unbounded and one is the single
point \mbox{$(n,i,j)=(10,1,2)$}.
In each region a selected exponent $w$ cannot equal any other
remaining exponent $v$. Writing $d=w-v$, the equality would require
\[
 d_A A+d_B B+d_C C=-d_0,\qquad A,B,C\in\mathbb Z_{\ge0}.
\]
All $768$ comparisons are refuted by $604$ sign contradictions,
$138$ gcd nondivisibilities and $26$ nonzero constants.
Thus $\Omega_3\ne0$ everywhere.

\subsection{The quintic collision cover}
The five-residue norm can be computed over $\mathbb F_2$.
If $P=p_0+\cdots+p_4$ is split by exponent residues modulo five,
the norm is the determinant of the circulant matrix
$(p_{(j-i)\bmod5})_{0\le i,j<5}$.
Its $26$ terms have the compact expression
\begin{equation}\label{eq:quintic-norm-form}
\begin{split}
 \Phi_5(p)={}&\sum_{i=0}^4p_i^5+\prod_{i=0}^4p_i\\
 &+\sum_{i=0}^4p_i^3
          (p_{i+1}p_{i-1}+p_{i+2}p_{i-2})\\
 &+\sum_{i=0}^4
       (p_i^2p_{i+1}^2p_{i+3}+p_i^2p_{i+2}^2p_{i+1}),
\end{split}
\end{equation}
with indices modulo five. Diagonalizing the circulant over
$\mathbb F_{16}$ gives the product in \eqref{eq:norm-product};
expansion of the $120$ determinant permutations independently
gives \eqref{eq:quintic-norm-form}.

The $125$ residue classes are initially divided into $467$ regions.
Of these, $354$ already have an uncancellable term and two have
identically zero defect. The remaining $111$ are resolved by the
following exact collision procedure.

Under the accumulated affine equalities, first group identical
exponents and cancel pairs. If the specialized polynomial is zero,
any selected remaining term must have a partner of equal exponent.
Branch over every feasible equality $E_w-E_v=0$. Each such branch
raises the coefficient rank: an equality of unchanged rank would
either have been cancelled already or be an inconsistent constant.
With three variables, the rank can rise at most to three.
The branches need only cover the zero set; they need not be disjoint.
An accepting node certifies inclusion in \eqref{eq:quintic-parameters};
it need not independently prove zero defect.
The explicit image formulas prove sufficiency separately.
Induction on the increasing rank therefore proves necessity:
each zero has a feasible partner at every nonterminal node,
and a unique-term or infeasible node cannot contain a zero.

Nonnegative integer feasibility is decided exactly.
At rank one, gcd divisibility is sufficient when coefficients have
mixed signs: an integer solution can be shifted into the nonnegative
orthant along a strictly positive integral kernel vector.
For coefficients of one sign, the target is either impossible by
sign or belongs to a finitely checked additive semigroup.
At rank two, Cramer's rule writes the variables as rational affine
functions of one original integer coordinate $t$.
Nonnegativity gives an interval, and the denominators give a finite
integrality period; an interval-and-residue test is necessary and
sufficient. The largest period in this certificate is three.
At rank three, test the unique rational point for nonnegative
integrality.

The resulting certificate has $672$ nodes: $540$ unique-term
nodes, $127$ collision nodes and five accepting nodes.
All $28\,611$ candidate-partner feasibility checks are verified.
The accepting sets are shown below; $A,B,C$ refer to the initial
gap variables, and $a,b,c$ to the positive gaps.
\begin{center}\small
\begin{tabular}{lll}
\toprule
$(k,i,j)\bmod5$ & conditions & resulting parameters\\
\midrule
$(0,3,4)$ & $A=B=C=0$ & $(10,3,4)$\\
$(3,1,2)$ & $A=B=C=0$ & $d=1$\\
$(1,2,1)$ & $A,B\ge1,\ C=0;\ a=2b-1,\ b=2c-1$ & $d=9$\\
$(1,2,1)$ & $A,B,C\ge1;\ a=2b-1,\ b=2c-1$ & $d=10t+9,\ t\ge1$\\
$(3,1,2)$ & $A,B,C\ge1;\ a=2b-1,\ b=2c-1$ & $d=10t+1,\ t\ge1$\\
\bottomrule
\end{tabular}
\end{center}
Putting $d=b$ in the last four rows gives
$c=(d+1)/2$, $a=2d-1$ and $n=7d-1$.
The residue conditions give exactly $d\equiv\pm1\pmod {10}$.
This proves the necessity part of Lemma~\ref{lem:norm-certificates}.

\begin{samepage}
\subsection{Independent verification and supplied files}
The following three commands replay the certificates in \path{verification/root_power/}.
\begin{quote}\small\ttfamily
python check\_cubic\_formal.py\\
python check\_quintic\_classification.py\\
python check\_quintic\_branches.py
\end{quote}
\end{samepage}
They use only the Python standard library.
The cubic checker reconstructs the norm from
$a^3+b^3+c^3+abc$, independently of the original
three-factor $\mathbb F_4$ construction.
The quintic classification checker uses the circulant determinant
and Cramer parametrization, while the construction uses
$\mathbb F_{16}$ products and Gaussian elimination.
It checks every partition, every feasible collision partner and
every accepting implication.
The branch checker separately verifies
\eqref{eq:quintic-support}, the ordered $22$-term supports,
the zero counts, the proper factors and the two small irreducible
inputs. It also uses $250$ polynomial substitutions to test the
implementation across all residue classes.

The complete finite certificates in that directory are
\begin{quote}\small\ttfamily
cubic\_certificate.json\\
quintic\_classification\_certificate.json
\end{quote}
The certificates and programs supply the finite arithmetic checks
for the parameter coverage proved above.
An additional independent audit, \path{check_certificate_lattice.py},
reconstructs both norms in the quotient algebra
$\mathbb F_2[y,z]/(z^\ell+y)$ and checks the same certificates using
unimodular integer solution lattices. It imports neither the generators
nor the three verification programs above.


\section{Proofs and certificates for the B family}\label{sec:B-certificates}

\subsection{Quadratic extension and the degree-\texorpdfstring{$12$}{12} period}
Write $A(x)=x^2+x+1$, $L_k(x)=\sum_{i=0}^{k-1}x^{2^i}$,
and $h_R^*(x)=x^{q+1}h_R(1/x)$. By Proposition~\ref{prop:B-descent},
\begin{equation}\label{eq:B-reciprocal}
 F_R(x):=x h_R^*(x)=A(x)x^q+(x+1)(L_k(x)+x^2)
                  =A(x)L_k(x)^2+x^2L_k(x)+x.
\end{equation}
The second expression uses $L_k(x)^2+L_k(x)=x^q+x$.
For $x\in\mathbb F_q$, this gives
$F_R(x)=x+(x+1)\operatorname{Tr}_{\mathbb F_q/\mathbb F_2}(x)$.
Its value is $x$ or $1$ according to the trace, so $h_R^*$ has no
roots in $\mathbb F_q$, since $h_R^*(0)=1$.

Let $x$ be a root of $h_R^*$ in $\mathbb F_{q^2}$ and put
$y=x^q$, $s=x+y$, $u=xy$, $v=(x+1)(y+1)$ and $t=L_k(x)$.
Then $s,v\ne0$, $s,u,v\in\mathbb F_q$ and $t^2+t=s$.
Solving \eqref{eq:B-reciprocal} for $t$ gives
\begin{align*}
 t &= xy+x^2+\frac{y}{x+1}=\alpha x+\beta,\qquad
 \alpha=s+\frac{s+1}{v},\qquad
 \beta=\frac{s^2+s+u}{v},\\
 c &:=\operatorname{Tr}_{\mathbb F_{q^2}/\mathbb F_2}(x)
       =t+t^q=\alpha s .
\end{align*}
To obtain the affine expression, use
$x^2=sx+u$ and $y(y+1)=(s+1)y+u$.
If $c=0$, then $v=(s+1)/s$, so $s\ne1$,
$u=(s+1)^2/s$ and $t=\beta=s^2+s+1$.
This gives $t^2+t=s^4+s$, contradicting $t^2+t=s$ and $s\ne0$.
Thus $c=1$. If $b$ were a root of $f_R$ in $\mathbb F_{q^2}$,
then $b\ne0,1$, since $f_R(0)=f_R(1)=1$. With
$a=b^2/(b+1)$, the element $a^{-1}$ is a root of $h_R^*$, and
$z=b/a$ satisfies $z^2+z=a^{-1}$. Taking absolute traces
contradicts the value $c=1$. This proves the first gcd in
Proposition~\ref{prop:B-quadratic-extension}.

For the other gcd, $c=1$ gives $\alpha=1/s$. If $s=1$ then
$v=u$, $t=x+1$, and $u=t^2+t=s=1$, so $A(x)=0$.
At either root of $A$, the values of $L_k(x)$ for $k\equiv0,1,2,3\pmod4$
are $0,x,1,x^2$. Equation~\eqref{eq:B-reciprocal} requires
$L_k(x)=x^2$, giving the stated quadratic factor.

If $s\ne1$, then $v=s/(s+1)$, $u=(s^2+s+1)/(s+1)$ and
$t=x/s+s^2+s+1/s$. Substituting into $t^2+t=s$ gives
$s^6+s/(s+1)=0$, or $s^6+s^5+1=0$. Write $G=X^6+X^5+1$.
The remainders of $X^4+X$ and $X^8+X$ modulo $G$ are
\texttt{0x12} and \texttt{0x25}, both coprime to $G$.
Thus $G$ has no irreducible factor of degree at most three and
is irreducible. It follows that $6\mid k$, that $u$ belongs to
$\mathbb F_{64}$, and that $x$ is quadratic over $\mathbb F_{64}$
as a root of $X^2+sX+u$. Here $x\notin\mathbb F_q$ and
$s\in\mathbb F_2(x)$, so the absolute degree of $x$ is $12$.
Consequently $k\equiv6\pmod{12}$ and $y=x^{64}$.
There are at most twelve candidate roots, from six possible values
of $s$ and two roots of $X^2+sX+u$ for each $s$.

The reciprocal $P_{12}^*$ has mask \texttt{0x18fb}. Modulo
$P_{12}^*$, $X^{2^{12}}=X$ and the remainders of
$X^{2^4}+X$, $X^{2^6}+X$ are \texttt{0x28b}, \texttt{0xaf3},
both coprime to the modulus. This proves irreducibility.
At $k=6$, the residues of $X^q$ and $L_6(X)$ are
\texttt{0xaf1} and \texttt{0x31b}, giving $F_7\equiv0$.
The twelve roots of $P_{12}^*$ therefore exhaust the candidates
of degree $12$ in the quadratic extension. Each has absolute
trace one by the trace argument at $k=6$.

To check divisibility of $h_R$ by $P_{12}$ for \emph{all} $k$, including
$6\nmid k$, take remainders of \eqref{eq:B-reciprocal} modulo
$P_{12}^*$. The term $X^{2^k}$ has period $12$, and
$L_{k+12}(X)+L_k(X)=1$ by the trace value. The following initial
remainders are hexadecimal, with $k=0,1$ used as residue representatives.
\begin{center}\small
\begin{tabular}{@{}r*{6}{c}@{}}
\toprule
$k$ & 0 & 1 & 2 & 3 & 4 & 5\\
Remainder & \texttt{0002} & \texttt{0016} & \texttt{0076} &
\texttt{0736} & \texttt{0e89} & \texttt{0731}\\
\midrule
$k$ & 6 & 7 & 8 & 9 & 10 & 11\\
Remainder & \texttt{0000} & \texttt{03ce} & \texttt{0c21} &
\texttt{0118} & \texttt{06a5} & \texttt{0d37}\\
\bottomrule
\end{tabular}
\end{center}
The next twelve remainders are obtained by adding $X+1$,
so $P_{12}$ occurs exactly at $k\equiv6\pmod{24}$.

The derivative of $F_R$ is $L_k^2+x^2+1$ by
\eqref{eq:B-reciprocal}. A common root of $F_R$ and its derivative
has $L_k(x)=x+1$, which makes $F_R(x)=A(x)^2$. Thus the degree-$12$
factor is simple. At the quadratic roots where $L_k(x)=x^2$,
the derivative vanishes, and the second Hasse derivative
$X^q+1$ has value $x\ne0$. This gives multiplicity two.
In each case the factor degree $e$ is $2$ or $12$ and $2k/e$
is odd, so the reciprocal root has trace one already in
$\mathbb F_{2^e}$. For a root $a$ of the corresponding factor of
$h_R$, the fibre $b^2+ab+a=0$ is therefore irreducible over
$\mathbb F_{2^e}$ by the trace criterion for $z^2+z=a^{-1}$.
Since $a=b^2/(b+1)$, the degree of $b$ is $2e$.
The fibres are separable and disjoint, so multiplicities are
preserved. This completes Proposition~\ref{prop:B-quadratic-extension}.

\subsection{Closed-orbit bound}
At any root $x$ of $h_R^*$, put $t=L_k(x)$.
Equation~\eqref{eq:B-reciprocal} gives $t\ne0,1$.
If $t\in\mathbb F_q$, then $x^q+x=t^2+t$ implies
$x\in\mathbb F_{q^2}$, where the preceding trace calculation
gives $t+t^q=1$, a contradiction. With $w=(t+1)/t$ we obtain
\begin{equation}\label{eq:B-orbit-relation}
 x^2+wx+\frac1w=0,\qquad x^q+x=\frac{w}{(w+1)^2}.
\end{equation}
For $x\in\mathbb F_{q^m}$, its conjugates
$x_i=x^{q^i}$ and $w_i=w^{q^i}$ satisfy
\begin{equation}\label{eq:B-orbit-system}
 w_i x_i^2+w_i^2x_i+1=0,\qquad
 (w_i+1)^2(x_{i+1}+x_i)+w_i=0
 \quad(i\in\mathbb Z/m\mathbb Z).
\end{equation}

The affine system has finitely many geometric solutions.
If it had a positive-dimensional component, take an irreducible
curve in that component and view the coordinates as rational
functions on its smooth projective model \cite[Section 53.2]{StacksCurves}.
Some $w_i$ is nonconstant, since if all $w_i$ were constant
the first equations would make every $x_i$ constant.
At a point where a nonconstant $w_i$ takes value $1$, the first
equation forces $x_i$ to be regular. A pole of $x_i$ would give
a unique term of highest pole order, namely $w_ix_i^2$.
Then the pole of $w_i/(w_i+1)^2$ forces $x_{i+1}$ to have a pole.
Whenever $x_j$ has a pole at this point, $w_j$ cannot take value
$1$. The function $W/(W+1)^2$ on the projective line is regular
away from $W=1$, including at infinity. Hence $x_{j+1}$ also has a pole.
Following the cycle forces $x_i$ to have a pole, a contradiction.

There are $2m$ equations of degree at most three in $2m$ variables,
so affine B\'ezout bounds their number of isolated solutions by
$9^m$ \cite[Theorem I.2]{MondalZeroes}. Distinct roots give
distinct first coordinates, and inversion transfers the bound
to $h_R$. Each root of $h_R$ gives at most two roots of $f_R$
in the same field. A factor of degree $d$ lies in the root count
as soon as $d\mid km$, whose least positive solution is
$m=d/\gcd(d,k)$. This proves Proposition~\ref{prop:B-closed-orbits}.

\subsection{Cubic extension and resultant certificates}
For the cubic classification we need the actual partial trace.
Put $s=t^2+t$, so $x^q=x+s$.
Adding the first equation in \eqref{eq:B-orbit-relation} to its
$q$th power and substituting $x^q=x+s$ gives
\begin{equation}\label{eq:B-trace-inverse}
 (w+w^q)x=s^2+w^q s+\frac1w+\frac1{w^q}.
\end{equation}
Since $t\notin\mathbb F_q$, the coefficient $w+w^q$ is nonzero.
Consequently $\mathbb F_2(t)=\mathbb F_2(x)$.

Suppose $24\mid k$ and $x$ is a root of $h_R^*$ in $\mathbb F_{q^3}$.
Its $q$-Frobenius orbit has length three by
Proposition~\ref{prop:B-quadratic-extension}. Write
$a=t$, $b=t^q$, $c=t^{q^2}$. These are distinct and avoid $0,1$.
Adding the three increments $x_{i+1}+x_i=t_i^2+t_i$
gives $c=a+b+\epsilon$ with $\epsilon\in\mathbb F_2$.
Define
\begin{align*}
 s&=a^2+a,\qquad u=b^2+b,\qquad D=(a+b)(a+1)(b+1),\\
 N&=\bigl(ab(s^2+s)+as\bigr)(a+1)(b+1)+ab(a+b),\qquad
 N_2=N+(s+u)D .
\end{align*}
By \eqref{eq:B-trace-inverse}, $x=N/D$ and $x^{q^2}=N_2/D$.
Let $Q(t,X)=(t^2+t)X^2+(t^2+1)X+t^2$.
For each $\epsilon$, take the resultant in $b$ of
$D^2Q(a,N/D)$ and $D^2Q(a+b+\epsilon,N_2/D)$,
after stripping all powers of $a,b,a+1,b+1,a+b,a+b+1$
from each polynomial. These six factors are nonzero on
the three-point orbit, so the stripping preserves all actual roots.

The two resultants have degree $128$ in $a$.
The following list gives the degree and multiplicity of each
irreducible factor. Repeated entries represent distinct factors.
\begin{center}\small
\begin{tabular}{@{}cl@{}}
\toprule
$\epsilon$ & (Degree, multiplicity)\\
\midrule
$0$ & $(1,7),(1,22),(3,3),(6,2),(14,1),(14,1),(14,1),(18,1),(18,1)$\\
$1$ & $(1,8),(1,24),(6,1),(6,2),(27,1),(51,1)$\\
\bottomrule
\end{tabular}
\end{center}
If $d$ is the degree of $a$, then $d/\gcd(d,k)=3$. As $24\mid k$,
this requires $9\mid d$. The resultant candidates are thus
the two degree-$18$ polynomials with masks \texttt{0x4fb89}
and \texttt{0x5c03d}, and the degree-$27$ polynomial
with mask \texttt{0xd01dfeb}.

For each candidate, let $a=X$ modulo its polynomial, set $b=a^{2^k}$,
and recover $x$ from \eqref{eq:B-trace-inverse}. This recovery
depends on $k\bmod d$, and $L_k(x)$ depends only on $k\bmod 2d$.
The following table covers every residue compatible with $24\mid k$
and $d/\gcd(d,k)=3$. Entries other than $k\bmod 2d$ are hexadecimal.
\begin{center}\small
\begin{tabular}{@{}lrrll@{}}
\toprule
Modulus & $k\bmod 2d$ & $x$ & $L_k(x)$ & $wx^2+w^2x+1$\\
\midrule
\texttt{4fb89} & 12 & \texttt{a8e2} & \texttt{11de9} & \texttt{2a7c0}\\
\texttt{4fb89} & 24 & \texttt{3588b} & \texttt{2} & \texttt{0}\\
\texttt{5c03d} & 12 & \texttt{29334} & \texttt{2} & \texttt{0}\\
\texttt{5c03d} & 24 & \texttt{22217} & \texttt{3c6cf} & \texttt{27b12}\\
\texttt{d01dfeb} & 18 & \texttt{fa84da} & \texttt{3} & \texttt{0}\\
\texttt{d01dfeb} & 36 & \texttt{77de2b8} & \texttt{1131750} & \texttt{72c5b6e}\\
\bottomrule
\end{tabular}
\end{center}
Both $L_k(x)=a$ and $wx^2+w^2x+1=0$ are required. They also
suffice by \eqref{eq:B-reciprocal} and \eqref{eq:B-orbit-relation}.
The degree-$27$ entry with zero final remainder has
$L_k(x)=a+1$, so it contributes no factor.
The two passing rows give the reciprocal factors of $P_{18,1}$
and $P_{18,2}$ by taking the Frobenius conjugates of $x$.
With $24\mid k$, their residues $24,12$ modulo $36$ give
$24,48$ modulo $72$, respectively.

The fibres above the degree-$18$ roots split in
$\mathbb F_{2^{36}}\subseteq\mathbb F_{q^3}$, since $36\mid3k$.
Thus the two corresponding gcds for $f_R$ are the entire pullbacks.
The mask of $\mathcal T(P_{18,1})$ is \texttt{0x105a535905},
with distinct irreducible factors of masks \texttt{0x52af3}
and \texttt{0x56ef3}. The mask of $\mathcal T(P_{18,2})$
is \texttt{0x1475b0a6d9} and it is irreducible.
The field polynomials $x^{q^2}+x$ and $x^{q^3}+x$ are squarefree,
so the root classifications determine the stated monic gcds.

\subsection{Replaying the certificates}
Run the following commands from \path{verification/B_family/}
using Python~3 with assertions enabled.
\begin{quote}\small\ttfamily
python -B verify\_quadratic\_extension.py\\
python -B verify\_cubic\_extension.py
\end{quote}
The quadratic checker uses integer bit arithmetic to verify the
small irreducibility tests, all $24$ residue classes for $P_{12}$,
and the pullback factors. The cubic checker computes and factors
both resultants over $\mathbb F_2$, tests the six candidate phases
and verifies the reciprocal factors and their pullbacks in the
passing cases. The cubic checker was tested with SymPy~1.14.0.

\clearpage

\begin{thebibliography}{99}
\small
\setlength{\itemsep}{1.5pt}
\setlength{\parsep}{0pt}
\setlength{\parskip}{0pt}
\bibitem{Coh06}
S. D. Cohen,
\emph{Primitive polynomials with a prescribed coefficient},
Finite Fields Appl. \textbf{12} (2006), 425--491.

\bibitem{LN}
R. Lidl and H. Niederreiter,
\emph{Finite Fields},
Encyclopedia of Mathematics and its Applications, vol.~20, 2nd ed., Cambridge University Press, 1997.

\bibitem{BC87}
J. V. Brawley and L. Carlitz,
\emph{Irreducibles and the composed product for polynomials over a finite field},
Discrete Math. \textbf{65} (1987), 115--139.
\href{https://doi.org/10.1016/0012-365X(87)90135-X}{doi:10.1016/0012-365X(87)90135-X}.

\bibitem{Swan62}
R. G. Swan,
\emph{Factorization of polynomials over finite fields},
Pacific J. Math. \textbf{12} (1962), 1099--1106.

\bibitem{FHS86}
H. M. Fredricksen, A. W. Hales, and M. M. Sweet,
\emph{A generalization of Swan's theorem},
Math. Comp. \textbf{46} (1986), 321--331.

\bibitem{Def20}
C. Defant,
\emph{Enumerating cliques in direct product graphs},
J. Combin. \textbf{11} (2020), 351--358.

\bibitem{Brouwer18}
A. E. Brouwer, S. M. Cioab\u{a}, F. Ihringer, and M. McGinnis,
\emph{The smallest eigenvalues of Hamming graphs, Johnson graphs and
other distance-regular graphs with classical parameters},
J. Combin. Theory Ser. B \textbf{133} (2018), 88--121.
\href{https://doi.org/10.1016/j.jctb.2018.04.005}{doi:10.1016/j.jctb.2018.04.005}.

\bibitem{Forney11}
G. D. Forney, Jr.,
\emph{Codes on graphs: duality and MacWilliams identities},
IEEE Trans. Inform. Theory \textbf{57} (2011), 1382--1397.
\href{https://doi.org/10.1109/TIT.2011.2104994}{doi:10.1109/TIT.2011.2104994}.

\bibitem{MP17}
S. Mattarei and M. Pizzato,
\emph{Generalizations of self-reciprocal polynomials},
Finite Fields Appl. \textbf{48} (2017), 271--288.
\href{https://doi.org/10.1016/j.ffa.2017.08.004}{doi:10.1016/j.ffa.2017.08.004}.

\bibitem{BFSS06}
A. Bostan, P. Flajolet, B. Salvy and \'E. Schost,
\emph{Fast computation of special resultants},
J. Symbolic Comput. \textbf{41} (2006), 1--29.
\href{https://doi.org/10.1016/j.jsc.2005.07.001}{doi:10.1016/j.jsc.2005.07.001}.

\bibitem{GHL15}
B. Grenet, J. van der Hoeven and G. Lecerf,
\emph{Randomized root finding over finite FFT-fields using tangent Graeffe transforms},
Proc. ISSAC 2015, 197--204.
\href{https://doi.org/10.1145/2755996.2756647}{doi:10.1145/2755996.2756647}.
Author version, Section 3.1:
\href{https://www.texmacs.org/joris/rmodroots/rmodroots.html}{generalized Graeffe transforms}.

\bibitem{BM25}
A. Bassa and R. Menares,
\emph{The $R$-transform as power map and its generalizations to higher degree},
Finite Fields Appl. \textbf{102} (2025), 102546.
\href{https://doi.org/10.1016/j.ffa.2024.102546}{doi:10.1016/j.ffa.2024.102546}.

\bibitem{Reis20}
L. Reis,
\emph{M\"obius-like maps on irreducible polynomials and rational transformations},
J. Pure Appl. Algebra \textbf{224} (2020), 169--180.
\href{https://doi.org/10.1016/j.jpaa.2019.05.001}{doi:10.1016/j.jpaa.2019.05.001}.

\bibitem{Long73}
A. F. Long, Jr.,
\emph{Factorization of irreducible polynomials over a finite field with
the substitution $x^{p^r}-x$ for $x$},
Duke Math. J. \textbf{40} (1973), 63--76.
\href{https://doi.org/10.1215/S0012-7094-73-04007-6}{doi:10.1215/S0012-7094-73-04007-6}.

\bibitem{Ago77}
S. Agou,
\emph{Factorisation sur un corps fini $\mathbb F_{p^n}$ des
polyn\^omes compos\'es $f(X^{p^r}-aX)$ lorsque $f(X)$ est un
polyn\^ome irr\'eductible de $\mathbb F_{p^n}(X)$},
J. Number Theory \textbf{9} (1977), 229--239.
\href{https://www.sciencedirect.com/science/article/pii/0022314X77900270}{Publisher page}.

\bibitem{Kap17}
G. Kapetanakis,
\emph{Prescribing coefficients of invariant irreducible polynomials},
J. Number Theory \textbf{180} (2017), 615--628.
\href{https://doi.org/10.1016/j.jnt.2017.05.003}{doi:10.1016/j.jnt.2017.05.003}.

\bibitem{Ugo13}
S. Ugolini,
\emph{Sequences of binary irreducible polynomials},
Discrete Math. \textbf{313} (2013), 2656--2662.
\href{https://doi.org/10.1016/j.disc.2013.08.011}{doi:10.1016/j.disc.2013.08.011}.

\bibitem{Elspas59}
B. Elspas,
\emph{The theory of autonomous linear sequential networks},
IRE Trans. Circuit Theory \textbf{6} (1959), 45--60.
\href{https://doi.org/10.1109/TCT.1959.1086506}{doi:10.1109/TCT.1959.1086506}.

\bibitem{Lempel70}
A. Lempel,
\emph{On a homomorphism of the de Bruijn graph and its applications
to the design of feedback shift registers},
IEEE Trans. Comput. \textbf{C-19} (1970), 1204--1209.
\href{https://doi.org/10.1109/T-C.1970.222859}{doi:10.1109/T-C.1970.222859}.

\bibitem{BZ04}
R. P. Brent and P. Zimmermann,
\emph{Algorithms for finding almost irreducible and almost primitive trinomials},
in Primes and Misdemeanours, Fields Inst. Commun., vol.~41,
Amer. Math. Soc., 2004, pp.~91--102.
\href{https://maths-people.anu.edu.au/~brent/pd/rpb212.pdf}{Author preprint}.

\bibitem{Flint}
The python-flint and FLINT developers,
\emph{python-flint}, version 0.9.0, with FLINT version 3.6.0.
\href{https://python-flint.readthedocs.io/}{Documentation};
\href{https://flintlib.org/doc/fmpz.html}{integer primality-testing documentation}.
Accessed 27 October 2026.

\bibitem{StacksCurves}
The Stacks Project Authors,
\emph{The Stacks Project}, Section 53.2, Curves and function fields.
\url{https://stacks.math.columbia.edu/tag/0BXX}.

\bibitem{MondalZeroes}
P. Mondal,
\emph{How many zeroes? Counting solutions of systems of polynomials
via toric geometry at infinity}, Theorem I.2.
\href{https://pinakimondal.org/books/p1_to_16.pdf}{Introduction and theorem}.

\end{thebibliography}
\end{document}
